\chapter{Generating Functions}\label{generating-functions}
\chapterstatus{complete}

\section{Introduction}\label{generating-functions:section-introduction}

A sequence of numbers $a_0, a_1, a_2, \ldots$ is an infinite object, and a counting problem usually defines it one term at a time. The idea of a \emph{generating function} is to package the
whole sequence into a single algebraic object, the formal power series $\sum_n a_nx^n$, and to translate operations on sequences into operations on series: adding sequences termwise is adding
series, the convolution $\sum_{i+j=n}a_ib_j$ is multiplication, the partial sums are division by $1 - x$, and a recurrence becomes an equation that the series satisfies. Solving that equation, by
the algebra of series, solves the counting problem. Wilf compares a generating function to a clothesline on which the sequence is hung up for display
\cite{wilf-generatingfunctionology}: all the terms can be seen, and handled, at once.

Here is the simplest example. The Fibonacci numbers $F_0 = 0$, $F_1 = 1$, $F_{n+2} = F_{n+1} + F_n$ have the generating function $F = \sum_nF_nx^n$, and the recurrence says that
$F - xF - x^2F = x$, so $F = x/(1 - x - x^2)$. Splitting this rational function into partial fractions gives at once the formula $F_n = (\varphi^n - \bar{\varphi}^n)/\sqrt 5$ with
$\varphi = (1 + \sqrt 5)/2$ and $\bar{\varphi} = (1 - \sqrt 5)/2$. Nothing here needs convergence: $x$ is a formal symbol, and every step is an identity in a ring of formal power series.

The chapter is organised as follows.
\begin{enumerate}
\item The ring $R[[x]]$ of formal power series is developed with complete proofs: units, the order of a series, infinite sums and products, composition and compositional inverses, the formal
derivative, and, over a ring containing $\QQ$, the exponential, the logarithm and the binomial series $(1 + x)^a$ for arbitrary exponents $a$
(Section~\ref{generating-functions:section-formal-power-series}).
\item The sequences satisfying a linear recurrence with constant coefficients are exactly those with a rational generating function, and partial fractions give their closed forms; Fibonacci and
Lucas numbers, words avoiding a pattern, and P\'olya's problem of making change are worked out (Section~\ref{generating-functions:section-rational}).
\item The Catalan numbers satisfy the quadratic equation $C = 1 + xC^2$; solving it gives $C_n = \frac{1}{n + 1}\binom{2n}{n}$, and bijections show that they count binary trees, plane trees, Dyck
paths, bracketings, triangulations and non-crossing partitions (Section~\ref{generating-functions:section-catalan}).
\item The generating function of the partitions of integers is an infinite product. From it we derive Euler's theorems again, the Gaussian binomial coefficients and the count of subspaces of a vector
space over a finite field, Euler's pentagonal number theorem with Franklin's involution, the recurrence for $p(n)$, and Jacobi's triple product identity with its corollaries
(Section~\ref{generating-functions:section-partitions}).
\item Exponential generating functions count labelled structures: products of exponential series count pairs of structures on complementary sets, and the exponential formula says that the
structures built from connected pieces have the exponential of the generating function of the pieces. Set partitions, permutations, involutions, derangements, connected graphs and Bernoulli numbers
illustrate it (Section~\ref{generating-functions:section-exponential}).
\item Lagrange's inversion formula computes the coefficients of a series defined implicitly by $T = x\phi(T)$; it gives Cayley's formula $n^{n-2}$ for the number of labelled trees, which we also
prove by Pitman's double counting, and the Fuss--Catalan numbers, which we also obtain from the cycle lemma (Section~\ref{generating-functions:section-lagrange}).
\item Elementary bounds for the coefficients, such as $4^n/(2\sqrt{n}) \leq \binom{2n}{n} \leq 4^n/\sqrt{2n + 1}$, conclude the chapter (Section~\ref{generating-functions:section-estimates}); the analytic theory
of the growth of coefficients needs complex analysis and is only described.
\end{enumerate}
References: \cite{wilf-generatingfunctionology}, \cite{stanley-ec1}, \cite[Chapters 5 and 6]{stanley-ec2}, \cite{flajolet-sedgewick} and, for partitions, \cite{andrews-partitions}.

\noindent\textbf{Prerequisites.} Chapter~\ref{rings} (Rings and Modules), Chapter~\ref{fields} (Fields and Galois Theory), Chapter~\ref{linear-algebra} (Linear Algebra) and
Chapter~\ref{combinatorics} (Combinatorics).

\noindent\textbf{On the way to the Hodge conjecture.} This chapter is not needed on the direct route; it belongs to the reading path of discrete mathematics
(Subsection~\ref{introduction:subsection-reading-paths}). Formal power series reappear as completions of polynomial rings in commutative algebra and as the local rings of smooth points in
algebraic geometry.

\begin{notation}\label{generating-functions:notation-basic}
Throughout, $R$ is a commutative ring with $1$ (Definition~\ref{rings:definition-ring}), and $K$ is a field. A ring \emph{contains $\QQ$} if it contains $\QQ$ as a subring with the same $1$;
equivalently, every nonzero integer is a unit in it. We use the notation of Chapter~\ref{combinatorics}: $[n] = \{1, \ldots, n\}$, $\binom{A}{k}$ is the set of $k$-element subsets of $A$, and
$x^{\underline{k}}$, $x^{\overline{k}}$ are the falling and rising factorials (Notation~\ref{combinatorics:notation-basic}). The binomial polynomials $\binom{x}{k} = x^{\underline{k}}/k!$
(Definition~\ref{combinatorics:definition-binomial-polynomial}) give $\binom{a}{k}$ for $a$ in any ring containing $\QQ$.
\end{notation}

\section{Formal power series}\label{generating-functions:section-formal-power-series}

A power series in analysis is a function, defined where the series converges. A \emph{formal} power series is just its sequence of coefficients, with the operations that the usual rules for
series would give; whether it converges anywhere is irrelevant. This section sets up the algebra carefully, because the rest of the chapter computes with it constantly, and because a few
operations that look harmless (substituting $1$ for $x$, taking $\exp$ of a series with nonzero constant term, multiplying infinitely many series) are meaningless and must be excluded.

\subsection{The ring of formal power series}\label{generating-functions:subsection-ring}

\begin{definition}\label{generating-functions:definition-power-series}
A \emph{formal power series} over $R$ is a sequence $f = (a_0, a_1, a_2, \ldots)$ of elements of $R$; the set of them is $R[[x]]$. On $R[[x]]$ we define the \emph{sum} $+$ and the \emph{product}
$\cdot$ by
\[
\begin{aligned}
&(a_n)_n + (b_n)_n = (a_n + b_n)_n, \qquad (a_n)_n \cdot (b_n)_n = (c_n)_n,\\
&\text{where}\quad c_n = \sum_{i + j = n}a_ib_j = \sum_{i=0}^{n}a_ib_{n-i}.
\end{aligned}
\]
The element $a_n$ is the \emph{coefficient of $x^n$} in $f$, written $[x^n]f$; the coefficient $a_0$ is the \emph{constant term}, also written $f(0)$. For a sequence $(a_n)_{n \geq 0}$ of
elements of $R$, the series $\sum_na_nx^n$ is its \emph{(ordinary) generating function}.
\end{definition}

\begin{notation}\label{generating-functions:notation-power-series}
The product is written by juxtaposition, $fg = f \cdot g$. We identify $r \in R$ with $(r, 0, 0, \ldots)$, write $x = (0, 1, 0, 0, \ldots)$, and, as Proposition~\ref{generating-functions:proposition-summable}
will justify, write the series $(a_n)_n$ as $\sum_{n \geq 0}a_nx^n$ or $a_0 + a_1x + a_2x^2 + \cdots$. A polynomial is a series with finitely many nonzero coefficients, so $R[x] \subseteq R[[x]]$
(Definition~\ref{rings:definition-polynomials}). When a series is called $f$, its coefficients are usually called $f_n$, so that $f = \sum_nf_nx^n$. The variable may have another name:
$R[[t]]$, $R[[q]]$, \ldots, are the same ring with the variable written $t$ or $q$.
\end{notation}

The product is the product of the familiar rule $(\sum a_ix^i)(\sum b_jx^j) = \sum_n(\sum_{i+j=n}a_ib_j)x^n$, obtained by multiplying out and collecting powers of $x$. The sum defining $c_n$ is finite,
so the product needs no convergence at all; this is the point of the definition.

\begin{lemma}\label{generating-functions:lemma-ring}
$R[[x]]$ is a commutative ring with zero $(0, 0, \ldots)$ and unit $1 = (1, 0, 0, \ldots)$, and $R[x]$ and $R$ are subrings. For $k \geq 0$, $x^k$ is the sequence with $1$ in position $k$ and $0$
elsewhere, and multiplication by $x^k$ shifts coefficients: $[x^n](x^kf) = [x^{n-k}]f$ for $n \geq k$, and $[x^n](x^kf) = 0$ for $n < k$.
\end{lemma}

\begin{proof}
\emph{Addition} is termwise, so it is associative and commutative, with neutral element $0$ and inverse $(-a_n)_n$.
\emph{Commutativity of the product}: $\sum_{i+j=n}a_ib_j = \sum_{j+i=n}b_ja_i$, since $R$ is commutative.
\emph{Associativity}: for series $f, g, h$ with coefficients $a_i, b_j, c_k$, the coefficient of $x^n$ in $(fg)h$ is $\sum_{m+k=n}(\sum_{i+j=m}a_ib_j)c_k = \sum_{i+j+k=n}a_ib_jc_k$, the sum over
the finitely many triples $(i, j, k) \in \NN^3$ with $i + j + k = n$; the same computation gives the same value for $f(gh)$.
\emph{Distributivity}: $\sum_{i+j=n}a_i(b_j + c_j) = \sum_{i+j=n}a_ib_j + \sum_{i+j=n}a_ic_j$.
\emph{Unit}: if $e = (1, 0, 0, \ldots)$ then $\sum_{i+j=n}e_ia_j = a_n$.
\emph{Subrings}: the sum and product of two sequences with finitely many nonzero terms again have finitely many nonzero terms, and the operations restricted to $R[x]$ are those of
Definition~\ref{rings:definition-polynomials}; the constant sequences $(r, 0, 0, \ldots)$ form a copy of $R$.
\emph{Shifts}: $x^k$ is the sequence with $1$ in position $k$ by Lemma~\ref{rings:lemma-polynomial-basics}, and $[x^n](x^kf) = \sum_{i+j=n}[x^i]x^k \cdot f_j = f_{n-k}$ if $n \geq k$, and $0$
otherwise.
\end{proof}

\begin{remark}\label{generating-functions:remark-formal}
A formal power series is not a function, and $x$ is not a number. The series $\sum_nn!\,x^n \in \ZZ[[x]]$ is a perfectly good element of $\ZZ[[x]]$, although, read as a real power series, it
converges only at $x = 0$. Conversely the real numbers $\sum_nx^n$ at $x = \frac12$ and at $x = -\frac12$ are irrelevant to the formal series $\sum_nx^n$. Identities between formal series are
identities between their coefficients, and they hold over every ring $R$ for which they make sense. When the series happen to converge, the formal identities give identities between functions; that
step uses analysis and is not taken in this chapter.
\end{remark}

\begin{definition}\label{generating-functions:definition-order}
The \emph{order} of a nonzero series $f$ is $\operatorname{ord}(f) = \min\{n : f_n \neq 0\}$, and $\operatorname{ord}(0) = \infty$, with the conventions $\infty + n = \infty$ and $n < \infty$
for $n \in \NN$.
\end{definition}

So $\operatorname{ord}(f) \geq k$ means that $f_0 = \cdots = f_{k-1} = 0$, that is, that $f = x^kg$ for a series $g$ (Lemma~\ref{generating-functions:lemma-ring}); the larger the order, the
``smaller'' the series. The order plays the role of the absolute value in analysis, upside down.

\begin{lemma}\label{generating-functions:lemma-order}
For $f, g \in R[[x]]$:
\begin{enumerate}
\item $\operatorname{ord}(f + g) \geq \min(\operatorname{ord}(f), \operatorname{ord}(g))$, with equality if $\operatorname{ord}(f) \neq \operatorname{ord}(g)$;
\item $\operatorname{ord}(fg) \geq \operatorname{ord}(f) + \operatorname{ord}(g)$, with equality if $R$ is an integral domain;
\item if $R$ is an integral domain, so is $R[[x]]$.
\end{enumerate}
\end{lemma}

\begin{proof}
(1) If $n < \min(\operatorname{ord} f, \operatorname{ord} g)$ then $f_n + g_n = 0 + 0 = 0$. If $\operatorname{ord} f = m < \operatorname{ord} g$, then $(f + g)_m = f_m + 0 \neq 0$.
(2) We may assume $f, g \neq 0$; let $m = \operatorname{ord} f$ and $k = \operatorname{ord} g$. A term $f_ig_j$ with $i + j < m + k$ has $i < m$ or $j < k$, so it vanishes; hence $(fg)_n = 0$ for
$n < m + k$. The coefficient of $x^{m+k}$ is $\sum_{i+j=m+k}f_ig_j = f_mg_k$, since the other terms vanish for the same reason; if $R$ is a domain, $f_mg_k \neq 0$.
(3) By (2), if $f, g \neq 0$ then $\operatorname{ord}(fg) < \infty$, that is, $fg \neq 0$.
\end{proof}

\begin{counterexample}\label{generating-functions:counterexample-order-zero-divisors}
Over a ring with zero divisors the inequality in (2) can be strict: in $(\ZZ/4\ZZ)[[x]]$, $(2 + x)(2 + 3x) = 4 + 8x + 3x^2 = 3x^2$, of order $2 > 0 + 0$. And $R[[x]]$ is never a domain when $R$
is not: $2 \cdot 2 = 0$ in $(\ZZ/4\ZZ)[[x]]$.
\end{counterexample}

The units of $R[[x]]$ were determined in Proposition~\ref{rings:proposition-power-series-units}: $f$ is invertible if and only if its constant term $f_0$ is invertible in $R$, and then the
coefficients of $g = f^{-1}$ are found one at a time from $g_0 = f_0^{-1}$ and $g_n = -f_0^{-1}\sum_{i=1}^{n}f_ig_{n-i}$. We write $1/f$ or $f^{-1}$ for the inverse, and $g/f$ for $gf^{-1}$. The
first examples are the most important.

\begin{example}[Geometric series]\label{generating-functions:example-geometric}
\begin{enumerate}
\item $(1 - x)(1 + x + x^2 + \cdots) = 1$, since the coefficient of $x^n$ in the product is $1 - 1 = 0$ for $n \geq 1$. So $\frac{1}{1 - x} = \sum_{n \geq 0}x^n$: the sequence $1, 1, 1, \ldots$
has the generating function $1/(1 - x)$.
\item More generally, for $a \in R$, $\frac{1}{1 - ax} = \sum_na^nx^n$, by the same computation: $(1 - ax)\sum_na^nx^n = \sum_na^nx^n - \sum_na^{n+1}x^{n+1} = 1$.
\item For $m \geq 1$,
\[
\frac{1}{(1 - x)^m} = \sum_{n \geq 0}\binom{n + m - 1}{m - 1}x^n.
\]
Indeed the coefficient of $x^n$ in the product of $m$ copies of $\sum_ex^e$ is the number of $m$-tuples $(e_1, \ldots, e_m) \in \NN^m$ with $e_1 + \cdots + e_m = n$, which is
$\binom{n + m - 1}{m - 1}$ by stars and bars (Theorem~\ref{combinatorics:theorem-stars-bars}). For $m = 2$: $\frac{1}{(1 - x)^2} = 1 + 2x + 3x^2 + \cdots = \sum_n(n + 1)x^n$.
\item Replacing $x$ by $ax$ (a substitution justified in Subsection~\ref{generating-functions:subsection-composition}, but here simply the same computation with $a^nx^n$ in place of $x^n$):
$\frac{1}{(1 - ax)^m} = \sum_n\binom{n + m - 1}{m - 1}a^nx^n$.
\end{enumerate}
\end{example}

\begin{counterexample}\label{generating-functions:counterexample-units}
The hypothesis on the constant term depends on the ring of coefficients. The series $2 - x$ is a unit in $\QQ[[x]]$, with inverse $\frac12 \cdot \frac{1}{1 - x/2} = \sum_n2^{-n-1}x^n$, but not in
$\ZZ[[x]]$: the constant term of $(2 - x)g$ is $2g_0 \neq 1$. And $x$ is not a unit in any $R[[x]]$ with $R \neq 0$: $xg$ has constant term $0$. In particular ``$1/x$'' is not a formal power series;
it is a formal Laurent series (Definition~\ref{generating-functions:definition-laurent}).
\end{counterexample}

\begin{example}[Partial sums]\label{generating-functions:example-partial-sums}
If $f = \sum_na_nx^n$, then $\frac{f}{1 - x} = f \cdot \sum_jx^j$ has coefficients $\sum_{i+j=n}a_i = a_0 + a_1 + \cdots + a_n$: dividing by $1 - x$ takes partial sums, and multiplying by
$1 - x$ takes differences $a_n - a_{n-1}$. For instance the numbers $0 + 1 + \cdots + n$ have the generating function $\frac{1}{1 - x}\sum_nnx^n$. Since $\sum_nnx^n = x\sum_n(n + 1)x^n =
\frac{x}{(1 - x)^2}$ by Example~\ref{generating-functions:example-geometric}(3), this is $\frac{x}{(1 - x)^3} = \sum_n\binom{n + 1}{2}x^n$, using (3) again with $m = 3$ and the shift of
Lemma~\ref{generating-functions:lemma-ring}. So $0 + 1 + \cdots + n = \binom{n + 1}{2}$, the formula that, according to the anecdote, Gauss found as a schoolboy.
\end{example}

Over a field the structure of $K[[x]]$ is very simple: every series is a power of $x$ times a unit.

\begin{proposition}\label{generating-functions:proposition-field-coefficients}
Let $K$ be a field.
\begin{enumerate}
\item Every nonzero $f \in K[[x]]$ is uniquely $f = x^mu$ with $m = \operatorname{ord}(f)$ and $u$ a unit of $K[[x]]$.
\item The ideals of $K[[x]]$ are $0$ and the ideals $(x^m)$, $m \geq 0$, which form the chain $K[[x]] = (1) \supsetneq (x) \supsetneq (x^2) \supsetneq \cdots$. In particular $K[[x]]$ is a principal
ideal domain with exactly one maximal ideal, $(x)$, and $(x)$ is the set of non-units.
\end{enumerate}
\end{proposition}

\begin{proof}
(1) Write $f = x^mu$ with $u_n = f_{n+m}$ (Lemma~\ref{generating-functions:lemma-ring}); then $u_0 = f_m \neq 0$ is invertible in $K$, so $u$ is a unit by
Proposition~\ref{rings:proposition-power-series-units}. If $x^mu = x^{m^{\prime}}u^{\prime}$ with $u, u^{\prime}$ units, comparing orders (Lemma~\ref{generating-functions:lemma-order}(2), as $K$ is a domain)
gives $m = m^{\prime}$, and cancelling $x^m$ in the domain $K[[x]]$ gives $u = u^{\prime}$.
(2) Let $I \neq 0$ be an ideal and $m$ the smallest order of a nonzero element of $I$, say of $f = x^mu$. Then $x^m = fu^{-1} \in I$, so $(x^m) \subseteq I$; and every $g \in I$ has order at least $m$,
hence $g = x^{\operatorname{ord} g}v \in (x^m)$. So $I = (x^m)$. These ideals are distinct, since $x^m \notin (x^{m+1})$ by comparing orders. A proper ideal is contained in $(x)$, so $(x)$ is the only
maximal ideal; a series lies in $(x)$ exactly when its constant term is $0$, that is, when it is not a unit.
\end{proof}

\begin{remark}\label{generating-functions:remark-several-variables}
Since $R$ is an arbitrary commutative ring, series in several variables need no new definition: $R[[x, y]] = R[[y]][[x]]$ consists of the families $(a_{m,n})$ indexed by $\NN^2$, written
$\sum a_{m,n}x^my^n$, and $R[y][[x]]$, the series in $x$ whose coefficients are polynomials in $y$, is a subring. Bivariate generating functions, which record two statistics at once (for example the
number of elements and the number of blocks of a set partition), live in such rings.
\end{remark}

\subsection{Infinite sums and products}\label{generating-functions:subsection-infinite-sums}

The expression $\sum_na_nx^n$ is an infinite sum of series $a_nx^n$, and many generating functions are infinite products, such as $\prod_{k \geq 1}(1 - x^k)^{-1}$ for partitions. Both make sense,
without any analysis, as long as each coefficient receives only finitely many contributions.

\begin{definition}\label{generating-functions:definition-summable}
A family $(f_i)_{i \in I}$ of series, indexed by an arbitrary set $I$, is \emph{summable} if for every $n$ there are only finitely many $i \in I$ with $[x^n]f_i \neq 0$. Its \emph{sum}
$\sum_{i \in I}f_i$ is the series whose coefficient of $x^n$ is the finite sum $\sum_{i \in I}[x^n]f_i$ (over the finitely many $i$ with $[x^n]f_i \neq 0$).
\end{definition}

A family $(f_i)$ is summable as soon as $\operatorname{ord}(f_i) \to \infty$ in the sense that for every $N$ only finitely many $i$ have $\operatorname{ord}(f_i) \leq N$: then $[x^n]f_i \neq 0$
forces $\operatorname{ord}(f_i) \leq n$. A finite family is always summable, with the usual sum. For example $(a_nx^n)_{n \in \NN}$ is summable, as $\operatorname{ord}(a_nx^n) \geq n$, and its sum
is the series $(a_n)_n$: this is what the notation $\sum_na_nx^n$ means.

\begin{proposition}\label{generating-functions:proposition-summable}
Let $(f_i)_{i \in I}$ and $(g_j)_{j \in J}$ be summable families.
\begin{enumerate}
\item (Linearity) For $r, s \in R$, the family $(rf_i)_{i \in I} \sqcup (sg_j)_{j \in J}$ (indexed by the disjoint union) is summable, with sum $r\sum_if_i + s\sum_jg_j$.
\item (Associativity) If $I = \bigsqcup_{\lambda \in \Lambda}I_\lambda$, then each $(f_i)_{i \in I_\lambda}$ is summable, the family of the sums $F_\lambda = \sum_{i \in I_\lambda}f_i$ is summable,
and $\sum_\lambda F_\lambda = \sum_if_i$. In particular the sum does not depend on the order of the terms.
\item (Distributivity) The family $(f_ig_j)_{(i, j) \in I \times J}$ is summable, and $\sum_{(i,j)}f_ig_j = (\sum_if_i)(\sum_jg_j)$.
\item (Series as sums) A series $f$ is the sum of the summable family $(f_nx^n)_{n \in \NN}$.
\end{enumerate}
\end{proposition}

\begin{proof}
Let $I(n) = \{i : [x^n]f_i \neq 0\}$, a finite set; similarly $J(n)$.
(1) The indices with a nonzero coefficient of $x^n$ lie in $I(n) \sqcup J(n)$, and the coefficient of the sum is $r\sum_{i \in I(n)}[x^n]f_i + s\sum_{j \in J(n)}[x^n]g_j$.
(2) The family $(f_i)_{i \in I_\lambda}$ is summable since $I_\lambda \cap I(n) \subseteq I(n)$ is finite. The coefficient $[x^n]F_\lambda$ is nonzero only if $I_\lambda$ meets $I(n)$, which happens for
finitely many $\lambda$ since $I(n)$ is finite and the $I_\lambda$ are disjoint; and $\sum_\lambda[x^n]F_\lambda = \sum_\lambda\sum_{i \in I_\lambda \cap I(n)}[x^n]f_i = \sum_{i \in I(n)}[x^n]f_i$ by
the associativity of finite sums (Lemma~\ref{combinatorics:lemma-finite-sums}).
(3) $[x^n](f_ig_j) = \sum_{p+q=n}[x^p]f_i[x^q]g_j$ is nonzero only if $i \in I(p)$ and $j \in J(q)$ for some $p + q = n$, which leaves finitely many pairs. Then
\[
\begin{aligned}
\sum_{(i,j)}[x^n](f_ig_j) &= \sum_{p+q=n}\Bigl(\sum_{i \in I(p)}[x^p]f_i\Bigr)\Bigl(\sum_{j \in J(q)}[x^q]g_j\Bigr)\\
&= \sum_{p+q=n}[x^p]\Bigl(\sum_if_i\Bigr)[x^q]\Bigl(\sum_jg_j\Bigr),
\end{aligned}
\]
which is the coefficient of $x^n$ in the product.
(4) $[x^m](f_nx^n)$ is $f_m$ for $n = m$ and $0$ otherwise.
\end{proof}

\begin{counterexample}\label{generating-functions:counterexample-not-summable}
The family $(1)_{n \in \NN}$ is not summable: ``$1 + 1 + 1 + \cdots$'' has no constant term. On the other hand $((-1)^nx^n)_{n \in \NN}$ is summable, with sum $\frac{1}{1 + x}$
(Example~\ref{generating-functions:example-geometric}(2)), although substituting $x = 1$ in it would give Grandi's series $1 - 1 + 1 - \cdots$, which has no sum. So one may never ``put $x = 1$'' in a
formal identity unless the series involved are polynomials. Summability is a condition on each coefficient separately, and whenever it holds,
Proposition~\ref{generating-functions:proposition-summable} allows any regrouping. This is unlike conditionally convergent series of real numbers, which can be rearranged to have any sum: a formal
sum never needs a limit of numbers, only finitely many additions per coefficient.
\end{counterexample}

\begin{definition}\label{generating-functions:definition-convergence}
A sequence $(f^{(k)})_{k \geq 0}$ of series \emph{converges} to $f$, written $f^{(k)} \to f$ or $f = \lim_kf^{(k)}$, if for every $n$ there is $k_0$ such that $[x^n]f^{(k)} = [x^n]f$ for all
$k \geq k_0$; equivalently, $\operatorname{ord}(f - f^{(k)}) \to \infty$.
\end{definition}

So convergence means that each coefficient is eventually constant. The partial sums $\sum_{n \leq N}a_nx^n$ converge to $\sum_na_nx^n$; more generally, if $(f_i)_{i \in \NN}$ is summable, then
$\sum_{i \leq N}f_i \to \sum_if_i$, since the coefficient of $x^n$ of the partial sum is final as soon as $N \geq \max I(n)$. Sums, products and, by
Proposition~\ref{generating-functions:proposition-summable}, summable sums of convergent sequences converge to what one expects: if $f^{(k)} \to f$ and $g^{(k)} \to g$, then
$f^{(k)}g^{(k)} \to fg$, because $[x^n](fg)$ involves only the coefficients of $x^0, \ldots, x^n$ of $f$ and $g$.

\begin{definition}\label{generating-functions:definition-infinite-product}
Let $(g_i)_{i \in I}$ be a family of series with $g_i(0) = 0$ for every $i$, such that for every $N$ only finitely many $i$ have $\operatorname{ord}(g_i) \leq N$. For a finite subset $S \subseteq I$
put $g_S = \prod_{i \in S}g_i$ (with $g_\emptyset = 1$). The \emph{infinite product} is
\[
\prod_{i \in I}(1 + g_i) = \sum_{S \subseteq I \text{ finite}}g_S.
\]
\end{definition}

\begin{proposition}\label{generating-functions:proposition-infinite-product}
Under the hypotheses of Definition~\ref{generating-functions:definition-infinite-product}:
\begin{enumerate}
\item the family $(g_S)_S$ is summable, so the product is defined, and its constant term is $1$;
\item if $I$ is finite, the infinite product is the usual product $\prod_{i \in I}(1 + g_i)$;
\item if $I = \NN$ (or any countable set enumerated as $i_0, i_1, \ldots$), the finite partial products $\prod_{k \leq N}(1 + g_{i_k})$ converge to $\prod_{i \in I}(1 + g_i)$;
\item if $I = I^{\prime} \sqcup I^{\prime\prime}$, then $\prod_I(1 + g_i) = \prod_{I^{\prime}}(1 + g_i) \cdot \prod_{I^{\prime\prime}}(1 + g_i)$; and the product of two such products,
$\prod_I(1 + g_i)\prod_J(1 + h_j)$, is the product over $I \sqcup J$.
\end{enumerate}
\end{proposition}

\begin{proof}
Let $F_n = \{i \in I : \operatorname{ord}(g_i) \leq n\}$, a finite set.
(1) By Lemma~\ref{generating-functions:lemma-order}(2), $\operatorname{ord}(g_S) \geq \sum_{i \in S}\operatorname{ord}(g_i)$. If $S$ contains an $i \notin F_n$, then $\operatorname{ord}(g_S) \geq \operatorname{ord}(g_i) > n$;
so $[x^n]g_S \neq 0$ forces $S \subseteq F_n$, and there are finitely many such $S$. The constant term is that of $g_\emptyset = 1$, since $\operatorname{ord}(g_S) \geq |S|$ (every $g_i$ has order at
least $1$).
(2) By induction on $|I|$: for $I = \emptyset$ both sides are $1$, and if $I = I_0 \sqcup \{j\}$, then $\prod_I(1 + g_i) = (\sum_{S \subseteq I_0}g_S)(1 + g_j) = \sum_{S \subseteq I_0}g_S +
\sum_{S \subseteq I_0}g_{S \cup \{j\}}$, which is the sum over all subsets of $I$, split according to whether they contain $j$.
(3) Let $I_N = \{i_0, \ldots, i_N\}$. By (2) the partial product is $\sum_{S \subseteq I_N}g_S$. Once $I_N \supseteq F_n$, the subsets $S$ with $[x^n]g_S \neq 0$ are all subsets of $I_N$, so the
coefficient of $x^n$ of the partial product equals that of the infinite product.
(4) The finite subsets of $I^{\prime} \sqcup I^{\prime\prime}$ are the unions $S^{\prime} \sqcup S^{\prime\prime}$, and $g_{S^{\prime} \sqcup S^{\prime\prime}} = g_{S^{\prime}}g_{S^{\prime\prime}}$; apply
Proposition~\ref{generating-functions:proposition-summable}(3). The second statement is the first for the family indexed by $I \sqcup J$, which satisfies the hypotheses of
Definition~\ref{generating-functions:definition-infinite-product}.
\end{proof}

\begin{example}\label{generating-functions:example-product-distinct}
The family $g_k = x^k$, $k \geq 1$, satisfies the hypotheses, and $\prod_{k \geq 1}(1 + x^k) = \sum_Sx^{\sum S}$, the sum over finite sets $S$ of positive integers. The coefficient of $x^n$ is the
number of ways to write $n$ as a sum of distinct positive integers, that is, the number of partitions of $n$ into distinct parts (Definition~\ref{combinatorics:definition-partition}):
\[
\prod_{k \geq 1}(1 + x^k) = 1 + x + x^2 + 2x^3 + 2x^4 + 3x^5 + 4x^6 + 5x^7 + \cdots.
\]
Similarly, with $g_k = \frac{x^k}{1 - x^k} = x^k + x^{2k} + \cdots$ (of order $k$), each factor $1 + g_k = \frac{1}{1 - x^k}$, and the infinite product $\prod_{k \geq 1}\frac{1}{1 - x^k}$ is defined;
Section~\ref{generating-functions:section-partitions} shows that it is the generating function of all partitions.
\end{example}

\begin{counterexample}\label{generating-functions:counterexample-products}
\begin{enumerate}
\item The product $\prod_{k \geq 1}(1 + x)$ of infinitely many copies of $1 + x$ is not defined: the coefficient of $x$ would be $1 + 1 + \cdots$. The condition that only finitely many $g_i$ have
order $\leq N$ excludes it.
\item The product $\prod_{k \geq 1}(1 + \frac{1}{2^k})$ of constants is not a formal product at all; it is a real number defined by a limit, and its value cannot be computed coefficientwise. The
condition $g_i(0) = 0$ excludes such factors. (A single factor with nonzero constant term, or finitely many, can of course be multiplied in afterwards.)
\end{enumerate}
\end{counterexample}

\begin{example}\label{generating-functions:example-binary-product}
$\prod_{k \geq 0}(1 + x^{2^k}) = \frac{1}{1 - x}$. Indeed the product is $\sum_Sx^{\sum_{k \in S}2^k}$ over finite sets $S \subseteq \NN$, and every $n \in \NN$ is the sum $\sum_{k \in S}2^k$ for
exactly one finite $S$, its binary expansion. So the infinite product of polynomials, none of which is a unit of finite order, is the geometric series: the uniqueness of binary expansions is a formal
identity. Equivalently $(1 - x)(1 + x)(1 + x^2)(1 + x^4) \cdots = 1$, which can be seen directly from $(1 - x)(1 + x) \cdots (1 + x^{2^{N}}) = 1 - x^{2^{N+1}}$ and
Proposition~\ref{generating-functions:proposition-infinite-product}(3).
\end{example}

\subsection{Composition}\label{generating-functions:subsection-composition}

In analysis one substitutes a function into a power series. Formally, one can substitute a series $g$ into $f = \sum_nf_nx^n$ as long as the family $(f_ng^n)$ is summable, which is guaranteed when
$g$ has no constant term.

\begin{definition}\label{generating-functions:definition-composition}
Let $f, g \in R[[x]]$ with $g(0) = 0$. The \emph{composition} of $f$ with $g$ is $f \circ g = f(g) = \sum_{n \geq 0}f_ng^n$, the sum of a summable family, since $\operatorname{ord}(f_ng^n) \geq n$
(Lemma~\ref{generating-functions:lemma-order}(2)).
\end{definition}

For instance $\frac{1}{1 - x} \circ ax = \sum_na^nx^n$, which is $\frac{1}{1 - ax}$ by Example~\ref{generating-functions:example-geometric}(2), and $\frac{1}{1 - x} \circ x^2 = \sum_nx^{2n}$. That
the composition with $g$ preserves all the algebra is the content of the next lemma.

\begin{lemma}\label{generating-functions:lemma-composition}
Let $g, h \in R[[x]]$ with $g(0) = h(0) = 0$.
\begin{enumerate}
\item The map $R[[x]] \to R[[x]]$, $f \mapsto f(g)$, is a ring homomorphism that is the identity on $R$ and sends $x$ to $g$. It preserves summable sums: if $(f_i)$ is summable, so is $(f_i(g))$,
and $(\sum_if_i)(g) = \sum_if_i(g)$. If $f$ is a unit, $(1/f)(g) = 1/f(g)$.
\item $\operatorname{ord}(f(g)) \geq \operatorname{ord}(f)$, and $f(g)(0) = f(0)$; so $f(g)(0) = 0$ when $f(0) = 0$, and then $(f \circ g) \circ h = f \circ (g \circ h)$.
\item $x \circ g = g$ and $f \circ x = f$.
\end{enumerate}
\end{lemma}

\begin{proof}
(1) Additivity is Proposition~\ref{generating-functions:proposition-summable}(1). For products, by Proposition~\ref{generating-functions:proposition-summable}(3) and (2),
\[
f(g)\,k(g) = \sum_{(i,j)}f_ik_jg^{i+j} = \sum_n\Bigl(\sum_{i+j=n}f_ik_j\Bigr)g^n = (fk)(g).
\]
Constants go to themselves and $x$ to $g$. For a summable family $(f_i)$, the coefficient $[x^n](f_i(g)) = \sum_{m \leq n}[x^m]f_i \cdot [x^n]g^m$ (the terms with $m > n$ vanish since
$\operatorname{ord}(g^m) \geq m$) is nonzero only if $[x^m]f_i \neq 0$ for some $m \leq n$, which happens for finitely many $i$; the family $(([x^m]f_i)\,g^m)_{(i,m)}$ is then summable, by the same argument, and regrouping it
in two ways (Proposition~\ref{generating-functions:proposition-summable}(2)) gives $\sum_if_i(g) = \sum_m(\sum_i[x^m]f_i)g^m = (\sum_if_i)(g)$. A ring homomorphism sends inverses to inverses.
(2) If $\operatorname{ord}(f) = m$, then $f(g) = \sum_{n \geq m}f_ng^n$ and each term has order at least $n \geq m$; the constant term of $f(g)$ is that of $f_0g^0 = f_0$. For associativity, both
$f \mapsto (f \circ g) \circ h$ and $f \mapsto f \circ (g \circ h)$ are ring homomorphisms (compositions of the homomorphisms of (1)) that preserve summable sums and send $x$ to $g \circ h$. Two such
maps agree on polynomials, being ring homomorphisms that agree on $R$ and on $x$, and hence on every $f = \sum_n f_nx^n$, the sum of the summable family $(f_nx^n)$.
(3) $x \circ g = g^1 = g$, and $f \circ x = \sum f_nx^n = f$.
\end{proof}

\begin{counterexample}\label{generating-functions:counterexample-composition}
The hypothesis $g(0) = 0$ cannot be dropped. Substituting $g = 1$ into $\frac{1}{1 - x} = \sum_nx^n$ would give $1 + 1 + \cdots$. Substituting $1 + x$ into $\sum_n\frac{x^n}{n!} \in \QQ[[x]]$ would give
a constant term $\sum_n\frac1{n!}$, a real number that is not a finite combination of the coefficients. When $f$ is a polynomial, every $g$ may be substituted, since the sum $\sum_nf_ng^n$ is
finite; for instance $(1 + x)^2 \circ (1 + x) = (2 + x)^2$.
\end{counterexample}

\begin{proposition}[Compositional inverse]\label{generating-functions:proposition-compositional-inverse}
Let $f \in R[[x]]$ with $f(0) = 0$. There is a series $g$ with $g(0) = 0$ and $f \circ g = x$ if and only if $f_1$ is a unit of $R$. In that case $g$ is unique, $g_1 = f_1^{-1}$, and also
$g \circ f = x$. The series with $f(0) = 0$ and $f_1$ a unit form a group under composition, with neutral element $x$.
\end{proposition}

\begin{proof}
\emph{Necessity.} The coefficient of $x$ in $f(g) = \sum_{n \geq 1}f_ng^n$ is $f_1g_1$, since $g^n$ has order at least $n \geq 2$ for $n \geq 2$. So $f \circ g = x$ forces $f_1g_1 = 1$.
\emph{Existence and uniqueness.} Suppose $f_1$ is a unit. For $n \geq 1$, the coefficient of $x^n$ in $f(g) = \sum_{m=1}^{n}f_mg^m + (\text{terms of order} > n)$ is
\[
[x^n]f(g) = f_1g_n + P_n(f_2, \ldots, f_n, g_1, \ldots, g_{n-1}),
\]
where $P_n$ is a polynomial with integer coefficients: in $g^m$ with $m \geq 2$, the coefficient of $x^n$ is a sum of products $g_{j_1} \cdots g_{j_m}$ with $j_1 + \cdots + j_m = n$ and all
$j_i \geq 1$, so all $j_i \leq n - 1$. The conditions $[x^1]f(g) = 1$ and $[x^n]f(g) = 0$ for $n \geq 2$ therefore determine $g_1 = f_1^{-1}$ and then recursively
$g_n = -f_1^{-1}P_n(f_2, \ldots, f_n, g_1, \ldots, g_{n-1})$: a solution exists and is unique.
\emph{The group.} Let $G$ be the set of series with $f(0) = 0$ and $f_1$ a unit. If $f, h \in G$, then $f \circ h \in G$, as $[x](f \circ h) = f_1h_1$ by the computation above; composition is associative
by Lemma~\ref{generating-functions:lemma-composition}(2), with neutral element $x$ by (3) of that lemma, and every $f \in G$ has a right inverse $g \in G$ ($g_1 = f_1^{-1}$ is a unit). In a monoid
in which every element has a right inverse, right inverses are two-sided: if $f \circ g = x$ and $g \circ h = x$, then $f = f \circ (g \circ h) = (f \circ g) \circ h = h$, so $g \circ f = g \circ h = x$.
\end{proof}

\begin{example}\label{generating-functions:example-compositional-inverse}
\begin{enumerate}
\item The inverse of $f = \frac{x}{1 - x}$ is $g = \frac{x}{1 + x}$: indeed $1 - g = \frac{1}{1 + x}$, so $f(g) = g \cdot (1 + x) = x$, using Lemma~\ref{generating-functions:lemma-composition}(1)
for $\frac{1}{1 - x} \circ g = \frac{1}{1 - g}$.
\item The inverse of $f = x - x^2$ is the series $g$ with $g - g^2 = x$. Its coefficients, computed from the recursion of the proof, begin $g = x + x^2 + 2x^3 + 5x^4 + 14x^5 + \cdots$, the Catalan
numbers; Section~\ref{generating-functions:section-catalan} solves this equation in closed form.
\end{enumerate}
\end{example}

\subsection{The formal derivative}\label{generating-functions:subsection-derivative}

\begin{definition}\label{generating-functions:definition-derivative}
The \emph{formal derivative} of $f = \sum_nf_nx^n$ is $Df = \sum_{n \geq 1}nf_nx^{n-1} = \sum_n(n + 1)f_{n+1}x^n$. We also write $f^{\prime}$ for $Df$, and $D^k$ for the $k$-th iterate.
\end{definition}

\begin{lemma}\label{generating-functions:lemma-derivative}
For $f, g \in R[[x]]$ and $r \in R$:
\begin{enumerate}
\item $D(f + g) = Df + Dg$, $D(rf) = rDf$, and $D$ preserves summable sums: $D\sum_if_i = \sum_iDf_i$;
\item (Leibniz rule) $D(fg) = (Df)g + fDg$; hence $D(f^n) = nf^{n-1}Df$ for $n \geq 1$, and $D(1/f) = -Df/f^2$ if $f$ is a unit;
\item (Chain rule) if $g(0) = 0$, then $D(f \circ g) = ((Df) \circ g) \cdot Dg$;
\item $D^kf$ has constant term $k!\,f_k$; so, if $R$ contains $\QQ$, $f_k = (D^kf)(0)/k!$ (Taylor's formula).
\end{enumerate}
\end{lemma}

\begin{proof}
(1) Linearity is clear coefficientwise; for a summable family, $[x^n]D\sum_if_i = (n + 1)\sum_i[x^{n+1}]f_i = \sum_i[x^n]Df_i$.
(2) Both sides of the Leibniz rule are additive in $f$ and in $g$ and preserve summable sums in each, by (1) and Proposition~\ref{generating-functions:proposition-summable}; writing $f$ and $g$ as sums of
monomials, it suffices to check it for $f = x^i$, $g = x^j$, where both sides are $(i + j)x^{i+j-1}$ (read as $0$ if $i = j = 0$). The formula for $D(f^n)$ follows by induction on $n$, and applying the
Leibniz rule to $f \cdot (1/f) = 1$ gives $(Df)/f + fD(1/f) = 0$.
(3) Both sides are additive in $f$ and preserve summable sums (Lemma~\ref{generating-functions:lemma-composition}(1) and (1)), so it suffices to take $f = x^n$: then $D(g^n) = ng^{n-1}Dg$ by (2),
and $(Dx^n) \circ g = ng^{n-1}$.
(4) $D^kf = \sum_n(n + 1)(n + 2) \cdots (n + k)f_{n+k}x^n$, by induction on $k$; its constant term is $k!\,f_k$.
\end{proof}

\begin{counterexample}\label{generating-functions:counterexample-characteristic}
In characteristic $p > 0$ the derivative loses information: in $\mathbb{F}_p[[x]]$, $D(x^p) = px^{p-1} = 0$, so a series with derivative $0$ need not be constant, and Taylor's formula is
meaningless since $p! = 0$. Over a ring containing $\QQ$ this cannot happen: if $Df = 0$ then $(n + 1)f_{n+1} = 0$ and hence $f_{n+1} = 0$ for every $n$, so $f = f_0$ is constant. Consequently,
over such a ring, \emph{a series is determined by its derivative and its constant term}, and every series $g$ has a unique antiderivative $\int g = \sum_n\frac{g_n}{n + 1}x^{n+1}$ with constant term
$0$.
\end{counterexample}

\subsection{Exponential, logarithm and powers}\label{generating-functions:subsection-exp-log}

From now on in this subsection, $R$ contains $\QQ$. The familiar series of analysis then make sense formally, and the familiar identities between them hold formally; but they must be proved again, by
algebra alone.

\begin{definition}\label{generating-functions:definition-exp-log}
The \emph{exponential} and \emph{logarithm} series are
\[
\begin{aligned}
\exp(x) &= e^x = \sum_{n \geq 0}\frac{x^n}{n!},\\
\log(1 + x) &= \sum_{n \geq 1}\frac{(-1)^{n-1}}{n}x^n = x - \frac{x^2}{2} + \frac{x^3}{3} - \cdots.
\end{aligned}
\]
For a series $f$ with $f(0) = 0$ we put $\exp(f) = \exp \circ f$ and $\log(1 + f) = \log(1 + x) \circ f$ (Definition~\ref{generating-functions:definition-composition}).
\end{definition}

The notation $\log(1 + f)$ is a single symbol: $\log$ is only defined for series with constant term $1$. Likewise $\exp(f)$ only for $f(0) = 0$
(Counterexample~\ref{generating-functions:counterexample-composition}).

\begin{proposition}\label{generating-functions:proposition-exp-log}
Let $f, g \in R[[x]]$ with $f(0) = g(0) = 0$.
\begin{enumerate}
\item $D\exp(x) = \exp(x)$, and $\exp(x)$ is the only series $E$ with $DE = E$ and $E(0) = 1$. Also $D\log(1 + x) = \frac{1}{1 + x}$.
\item $\exp(f + g) = \exp(f)\exp(g)$; in particular $\exp(f)$ is a unit with inverse $\exp(-f)$, and $\exp(nf) = \exp(f)^n$ for $n \in \ZZ$.
\item $\exp(\log(1 + f)) = 1 + f$ and $\log(\exp(f)) = f$, where $\log(\exp(f))$ means $\log(1 + (\exp(f) - 1))$.
\item $\log((1 + f)(1 + g)) = \log(1 + f) + \log(1 + g)$.
\item $D\exp(f) = \exp(f)Df$ and $D\log(1 + f) = \frac{Df}{1 + f}$.
\end{enumerate}
\end{proposition}

\begin{proof}
(1) $D\sum_n\frac{x^n}{n!} = \sum_{n \geq 1}\frac{nx^{n-1}}{n!} = \sum_m\frac{x^m}{m!}$. If $DE = E$ and $E(0) = 1$, then $(n + 1)E_{n+1} = E_n$, so $E_n = 1/n!$ by induction. Finally
$D\log(1 + x) = \sum_{n \geq 1}(-1)^{n-1}x^{n-1} = \frac{1}{1 + x}$ by Example~\ref{generating-functions:example-geometric}(2) with $a = -1$.
(2) By Proposition~\ref{generating-functions:proposition-summable}(3) and (2), and the binomial theorem (Theorem~\ref{combinatorics:theorem-binomial}) in the commutative ring $R[[x]]$,
\[
\begin{aligned}
\exp(f)\exp(g) &= \sum_{(i,j)}\frac{f^ig^j}{i!\,j!} = \sum_n\frac{1}{n!}\sum_{i+j=n}\binom{n}{i}f^ig^j\\
&= \sum_n\frac{(f + g)^n}{n!} = \exp(f + g);
\end{aligned}
\]
the family $(f^ig^j/(i!j!))$ is summable since its terms have order at least $i + j$. With $g = -f$ we get $\exp(f)\exp(-f) = \exp(0) = 1$, and $\exp(nf) = \exp(f)^n$ follows by induction for
$n \geq 0$ and then for $n < 0$ by taking inverses.
(5) Chain rule (Lemma~\ref{generating-functions:lemma-derivative}(3)) and (1): $D\exp(f) = \exp(f)Df$, and $D\log(1 + f) = \frac{1}{1 + x} \circ f \cdot Df = \frac{Df}{1 + f}$, using
Lemma~\ref{generating-functions:lemma-composition}(1) for the inverse.
(3) Let $L = \log(1 + x)$ and $h = \exp(L)$. By (5), $Dh = h \cdot \frac{1}{1 + x}$, so by the Leibniz rule
\[
D\Bigl(\frac{h}{1 + x}\Bigr) = \frac{Dh}{1 + x} - \frac{h}{(1 + x)^2} = 0.
\]
By Counterexample~\ref{generating-functions:counterexample-characteristic} the series $h/(1 + x)$ is constant, equal to its constant term $h(0) = \exp(0) = 1$; so $\exp(\log(1 + x)) = 1 + x$. Put
$E_1 = \exp(x) - 1$. We have shown $E_1 \circ L = x$. Both $E_1$ and $L$ have constant term $0$ and coefficient of $x$ equal to $1$, so they lie in the group of
Proposition~\ref{generating-functions:proposition-compositional-inverse}, where $E_1 \circ L = x$ implies $L \circ E_1 = x$, that is, $\log(\exp(x)) = x$. Composing both identities with $f$
(Lemma~\ref{generating-functions:lemma-composition}(2)) gives (3).
(4) By (2) and (3), $\exp(\log(1 + f) + \log(1 + g)) = (1 + f)(1 + g) = 1 + (f + g + fg)$. Apply $\log$, using $\log(\exp(u)) = u$ for $u = \log(1 + f) + \log(1 + g)$, which has constant term $0$.
\end{proof}

\begin{definition}\label{generating-functions:definition-powers}
For $a \in R$ and a series $f$ with $f(0) = 0$, put $(1 + f)^a = \exp(a\log(1 + f))$. The \emph{binomial series} is $(1 + x)^a$.
\end{definition}

\begin{theorem}[Newton's binomial series]\label{generating-functions:theorem-binomial-series}
Let $R$ contain $\QQ$, $a, b \in R$, and $f, g \in R[[x]]$ with $f(0) = g(0) = 0$.
\begin{enumerate}
\item $(1 + x)^a = \sum_{n \geq 0}\binom{a}{n}x^n$, with $\binom{a}{n} = \frac{a(a - 1) \cdots (a - n + 1)}{n!}$.
\item $(1 + f)^a(1 + f)^b = (1 + f)^{a+b}$, $((1 + f)^a)^b = (1 + f)^{ab}$, and $((1 + f)(1 + g))^a = (1 + f)^a(1 + g)^a$; here $((1 + f)^a)^b$ makes sense since $(1 + f)^a$ has constant term $1$.
\item For $a = n \in \NN$, $(1 + f)^a$ is the $n$-th power $(1 + f)^n$ in the ring $R[[x]]$; for $a = -1$ it is the inverse $1/(1 + f)$.
\item For $m \geq 1$, $(1 + f)^{1/m}$ is the unique series $h$ with constant term $1$ and $h^m = 1 + f$.
\end{enumerate}
\end{theorem}

\begin{proof}
(1) Let $B = (1 + x)^a = \exp(aL)$ with $L = \log(1 + x)$. By Proposition~\ref{generating-functions:proposition-exp-log}(5) and (1), $DB = B \cdot \frac{a}{1 + x}$, that is, $(1 + x)DB = aB$.
Comparing coefficients of $x^n$: $(n + 1)B_{n+1} + nB_n = aB_n$, so $B_{n+1} = \frac{a - n}{n + 1}B_n$, with $B_0 = \exp(0) = 1$. By induction $B_n = \frac{a(a - 1) \cdots (a - n + 1)}{n!} = \binom{a}{n}$.
(2) The first two identities: $\exp(a\log(1 + f))\exp(b\log(1 + f)) = \exp((a + b)\log(1 + f))$ by Proposition~\ref{generating-functions:proposition-exp-log}(2); and
$\log((1 + f)^a) = \log(\exp(a\log(1 + f))) = a\log(1 + f)$ by (3) of that proposition, so $((1 + f)^a)^b = \exp(ab\log(1 + f))$. The third: $\log((1 + f)(1 + g)) = \log(1 + f) + \log(1 + g)$
by (4) of that proposition; multiply by $a$ and exponentiate.
(3) By (2), $(1 + f)^{n} = \exp(\log(1 + f))^n$ (using $\exp(nu) = \exp(u)^n$), and $\exp(\log(1 + f)) = 1 + f$. For $a = -1$, $(1 + f)^{-1}(1 + f)^1 = (1 + f)^0 = 1$.
(4) By (2) and (3), $h = (1 + f)^{1/m}$ satisfies $h^m = (1 + f)^{1} = 1 + f$. If also $k^m = 1 + f$ with $k(0) = 1$, then $u = k/h$ has $u^m = 1$ and $u(0) = 1$; write $u = 1 + v$ with
$v(0) = 0$. If $v \neq 0$, of order $r \geq 1$, then $u^m = 1 + mv + (\text{terms of order} \geq 2r)$, so the coefficient of $x^r$ of $u^m - 1$ is $mv_r$, which is nonzero because $m$ is a unit of $R$ and $v_r \neq 0$; a contradiction. So $v = 0$ and $k = h$.
\end{proof}

\begin{example}\label{generating-functions:example-central-binomial}
For $a = -\frac12$, $\binom{-1/2}{n} = \frac{(-\frac12)(-\frac32) \cdots (-\frac{2n-1}{2})}{n!} = \frac{(-1)^n\,1 \cdot 3 \cdots (2n - 1)}{2^nn!} = \frac{(-1)^n(2n)!}{4^n\,n!^2}$, using
$1 \cdot 3 \cdots (2n - 1) = \frac{(2n)!}{2 \cdot 4 \cdots 2n} = \frac{(2n)!}{2^nn!}$. Substituting $-4x$ for $x$ (Lemma~\ref{generating-functions:lemma-composition}):
\[
\begin{aligned}
\frac{1}{\sqrt{1 - 4x}} = (1 - 4x)^{-1/2} &= \sum_{n \geq 0}\binom{2n}{n}x^n\\
&= 1 + 2x + 6x^2 + 20x^3 + 70x^4 + \cdots.
\end{aligned}
\]
Squaring and using Theorem~\ref{generating-functions:theorem-binomial-series}(2), $\bigl(\sum_n\binom{2n}{n}x^n\bigr)^2 = (1 - 4x)^{-1} = \sum_n4^nx^n$, which is the identity
\[
\sum_{k=0}^{n}\binom{2k}{k}\binom{2n - 2k}{n - k} = 4^n.
\]
Combinatorial proofs of this identity exist, but they are surprisingly delicate; the generating function proves it in one line.
\end{example}

\begin{example}\label{generating-functions:example-sqrt}
Similarly $\binom{1/2}{n} = \frac{1}{1 - 2n}\binom{-1/2}{n}$: the numerator of $\binom{1/2}{n}$ is $\frac12(\frac12 - 1) \cdots (\frac12 - n + 1)$, and that of $\binom{-1/2}{n}$ is
$(\frac12 - 1) \cdots (\frac12 - n + 1)(\frac12 - n)$, so their ratio is $\frac{1/2}{1/2 - n} = \frac{1}{1 - 2n}$. Hence
\[
\begin{aligned}
\sqrt{1 - 4x} = (1 - 4x)^{1/2} &= \sum_{n \geq 0}\frac{1}{1 - 2n}\binom{2n}{n}x^n\\
&= 1 - 2x - 2x^2 - 4x^3 - 10x^4 - 28x^5 - \cdots,
\end{aligned}
\]
and for $n \geq 1$ the coefficient $-\frac{1}{2n - 1}\binom{2n}{n} = -\frac{2}{n}\binom{2n - 2}{n - 1}$ is $-2$ times a Catalan number $\frac{1}{n}\binom{2n-2}{n-1}$; the check:
$\binom{2n}{n} = \frac{2n(2n - 1)}{n^2}\binom{2n - 2}{n - 1}$.
\end{example}

\begin{remark}\label{generating-functions:remark-formal-analytic}
When $R = \RR$ or $\CC$ and the series converge, the formal $\exp$, $\log$ and $(1 + x)^a$ are the Taylor series of the functions of the same name, and the formal identities become identities of
functions near $0$. The formal theory needs none of this, and it works over rings such as $\QQ[y]$ or $\QQ[[y]]$, where $a$ may itself be a variable: $(1 + x)^y = \sum_n\binom{y}{n}x^n$ in
$\QQ[y][[x]]$.
\end{remark}

\section{Rational generating functions and linear recurrences}\label{generating-functions:section-rational}

The simplest sequences after the polynomial ones are those defined by a linear recurrence with constant coefficients, such as the Fibonacci numbers. Their generating functions are exactly the
rational functions, and the factorisation of the denominator gives a closed formula.

\subsection{Recurrences and rational series}\label{generating-functions:subsection-recurrences}

\begin{definition}\label{generating-functions:definition-linear-recurrence}
Let $c_1, \ldots, c_d \in R$, $d \geq 1$. A sequence $(a_n)_{n \geq 0}$ in $R$ satisfies the \emph{linear recurrence with constant coefficients}
\[
a_{n+d} = c_1a_{n+d-1} + c_2a_{n+d-2} + \cdots + c_da_n \qquad (n \geq 0)
\]
of order $d$. Its \emph{characteristic polynomial} is $\chi(t) = t^d - c_1t^{d-1} - \cdots - c_d \in R[t]$, and its \emph{reflected polynomial} is $Q(x) = 1 - c_1x - \cdots - c_dx^d \in R[x]$.
\end{definition}

A sequence satisfying such a recurrence is determined by its $d$ initial values $a_0, \ldots, a_{d-1}$, and any initial values can be chosen.

\begin{theorem}\label{generating-functions:theorem-rational}
Let $c_1, \ldots, c_d \in R$ and $Q = 1 - c_1x - \cdots - c_dx^d$. For a sequence $(a_n)$ in $R$ with generating function $A = \sum_na_nx^n$, the following are equivalent:
\begin{enumerate}
\item $(a_n)$ satisfies the recurrence $a_{n+d} = c_1a_{n+d-1} + \cdots + c_da_n$ for all $n \geq 0$;
\item $QA$ is a polynomial of degree less than $d$ (or zero);
\item $A = P/Q$ for a polynomial $P$ of degree less than $d$ (or $P = 0$).
\end{enumerate}
In that case $P = QA$ is determined by $a_0, \ldots, a_{d-1}$: its coefficient of $x^m$ is $a_m - c_1a_{m-1} - \cdots - c_ma_0$ for $m < d$. Conversely, for every $P$ of degree less than $d$, the
coefficients of $P/Q$ satisfy the recurrence.
\end{theorem}

\begin{proof}
$Q$ has constant term $1$, so it is a unit (Proposition~\ref{rings:proposition-power-series-units}) and (2) and (3) are equivalent. Put $c_0 = -1$, so that $Q = -\sum_{k=0}^{d}c_kx^k$. For
$m \geq 0$, the coefficient of $x^m$ in $QA$ is $-\sum_{k=0}^{\min(m, d)}c_ka_{m-k}$. For $m = n + d \geq d$ this is $a_{n+d} - c_1a_{n+d-1} - \cdots - c_da_n$, which vanishes for all $n$ exactly when
(1) holds; and for $m < d$ it is the stated coefficient of $P$.
\end{proof}

\begin{example}\label{generating-functions:example-fibonacci-gf}
The Fibonacci numbers $F_0 = 0$, $F_1 = 1$, $F_{n+2} = F_{n+1} + F_n$ satisfy the recurrence with $d = 2$, $c_1 = c_2 = 1$, so $Q = 1 - x - x^2$, and $P = F_0 + (F_1 - F_0)x = x$:
\[
\sum_nF_nx^n = \frac{x}{1 - x - x^2}.
\]
The \emph{Lucas numbers} $L_0 = 2$, $L_1 = 1$, $L_{n+2} = L_{n+1} + L_n$, that is $2, 1, 3, 4, 7, 11, 18, 29, \ldots$, satisfy the same recurrence with $P = 2 + (1 - 2)x = 2 - x$. Both are sequences of
integers; the theorem holds over $\ZZ$.
\end{example}

\begin{counterexample}\label{generating-functions:counterexample-variable-coefficients}
The hypothesis that the coefficients are constant is essential. The sequence $a_n = n!$ satisfies the recurrence $a_{n+1} = (n + 1)a_n$, whose coefficient depends on $n$, but $\sum_nn!\,x^n$ is not
rational, even over $\CC$. Indeed, by Theorem~\ref{generating-functions:theorem-closed-form} below, the coefficients of a rational series $P/Q \in \CC[[x]]$ satisfy $|a_n| \leq Cn^dM^n$ for some
constants $C, M$ and $d$, while $n! \geq (n/2)^{n/2}$ (the $\lceil n/2\rceil$ largest factors are at least $n/2$), and $(n/2)^{n/2}/(n^dM^n) = (\sqrt{n/2}/M)^n/n^d$ is unbounded.
\end{counterexample}

\subsection{Partial fractions and the closed formula}\label{generating-functions:subsection-partial-fractions}

\begin{lemma}[Partial fractions]\label{generating-functions:lemma-partial-fractions}
Let $K$ be a field, and let $Q = \prod_{i=1}^{r}(1 - \alpha_ix)^{m_i}$ with $\alpha_1, \ldots, \alpha_r \in K$ distinct and nonzero and $m_i \geq 1$; put $d = \deg Q = m_1 + \cdots + m_r$. For every
polynomial $P \in K[x]$ of degree less than $d$ there are unique $\beta_{ij} \in K$ ($1 \leq i \leq r$, $1 \leq j \leq m_i$) with
\[
\frac{P}{Q} = \sum_{i=1}^{r}\sum_{j=1}^{m_i}\frac{\beta_{ij}}{(1 - \alpha_ix)^j}.
\]
\end{lemma}

\begin{proof}
\emph{Step 1: the case $r = 1$.} Write $\alpha = \alpha_1$, $m = m_1$ and $y = 1 - \alpha x$. Since $\alpha \neq 0$, $x = \alpha^{-1}(1 - y)$, so every polynomial of degree less than $m$ in $x$ is a
polynomial of degree less than $m$ in $y$: $P = \sum_{k=0}^{m-1}\gamma_ky^k$. Dividing by $y^m$ gives $P/Q = \sum_{k=0}^{m-1}\gamma_ky^{k-m}$, the required form with $\beta_{1j} = \gamma_{m-j}$.
\emph{Step 2: splitting.} Let $Q_i = (1 - \alpha_ix)^{m_i}$. For $i \neq k$ the polynomials $Q_i$ and $Q_k$ have no common root (the only root of $Q_i$ is $\alpha_i^{-1}$), so $Q_1$ and
$Q_2 \cdots Q_r$ generate the unit ideal of the principal ideal domain $K[x]$ (Theorem~\ref{rings:theorem-polynomial-division}): a generator of the ideal they generate divides both, and a common
divisor of positive degree would have a root, in an extension of $K$, common to $Q_1$ and $Q_2 \cdots Q_r$. Write $1 = UQ_1 + VQ_2 \cdots Q_r$ with $U, V \in K[x]$. Then
$\frac{P}{Q} = \frac{PV}{Q_1} + \frac{PU}{Q_2 \cdots Q_r}$. Divide: $PV = SQ_1 + P_1$ with $\deg P_1 < m_1$. Then $\frac{P}{Q} = \frac{P_1}{Q_1} + \frac{P_2}{Q_2 \cdots Q_r}$ with
$P_2 = PU + SQ_2 \cdots Q_r$; multiplying by $Q$ shows $P_2Q_1 = P - P_1Q_2 \cdots Q_r$, which has degree less than $d$, so $\deg P_2 < d - m_1$. By induction on $r$ and Step 1, the form exists.
\emph{Step 3: uniqueness.} Multiplying by $Q$, the right-hand side becomes $\sum_{i,j}\beta_{ij}Q/(1 - \alpha_ix)^j$, a polynomial of degree less than $d$. So the $K$-linear map
$K^d \to \{P : \deg P < d\}$, $(\beta_{ij}) \mapsto \sum_{i,j}\beta_{ij}Q/(1 - \alpha_ix)^j$, is surjective by Step 2; both spaces have dimension $d$, so it is injective
(Theorem~\ref{linear-algebra:theorem-rank-nullity}).
\end{proof}

\begin{theorem}\label{generating-functions:theorem-closed-form}
Let $K$ be a field, $c_1, \ldots, c_d \in K$ with $c_d \neq 0$, and suppose that the characteristic polynomial factors as $\chi(t) = \prod_{i=1}^{r}(t - \alpha_i)^{m_i}$ with distinct
$\alpha_i \in K$ (this holds for $K = \CC$ by the fundamental theorem of algebra, Theorem~\ref{fields:theorem-fta}). Then the sequences satisfying the recurrence are exactly the sequences
\[
a_n = \sum_{i=1}^{r}\sum_{j=1}^{m_i}\beta_{ij}\binom{n + j - 1}{j - 1}\alpha_i^n \qquad (\beta_{ij} \in K).
\]
If $K$ has characteristic $0$, they are exactly the sequences $a_n = \sum_ip_i(n)\alpha_i^n$ with polynomials $p_i \in K[t]$ of degree less than $m_i$.
\end{theorem}

\begin{proof}
\emph{The reflected polynomial.} The roots $\alpha_i$ are nonzero, since $\chi(0) = -c_d \neq 0$. For a polynomial $f$ of degree at most $e$ put $f^{\ast} = x^ef(1/x)$, that is,
$(\sum_kf_kt^k)^{\ast} = \sum_kf_kx^{e-k}$; then $(fg)^{\ast} = f^{\ast}g^{\ast}$ when the degrees are bounded by $e$ and $e^{\prime}$ and $fg$ is read with the bound $e + e^{\prime}$, as one checks on
monomials. So $Q = \chi^{\ast} = \prod_i((t - \alpha_i)^{\ast})^{m_i} = \prod_i(1 - \alpha_ix)^{m_i}$.
\emph{The formula.} By Theorem~\ref{generating-functions:theorem-rational}, the solutions are the coefficient sequences of the series $P/Q$ with $\deg P < d$, and by
Lemma~\ref{generating-functions:lemma-partial-fractions} these are the series $\sum_{i,j}\beta_{ij}(1 - \alpha_ix)^{-j}$ with arbitrary $\beta_{ij}$. By
Example~\ref{generating-functions:example-geometric}(4), $[x^n](1 - \alpha x)^{-j} = \binom{n + j - 1}{j - 1}\alpha^n$.
\emph{Characteristic $0$.} Here $\binom{n + j - 1}{j - 1}$ is the value at $n$ of the polynomial $\binom{t + j - 1}{j - 1} = (t + 1) \cdots (t + j - 1)/(j - 1)!$ of degree $j - 1$, and the polynomials
$\binom{t + j - 1}{j - 1}$, $1 \leq j \leq m$, form a basis of the polynomials of degree less than $m$, since their degrees are $0, 1, \ldots, m - 1$. So the combinations
$\sum_j\beta_{ij}\binom{n + j - 1}{j - 1}$ are exactly the values $p_i(n)$ of the polynomials of degree less than $m_i$. (In characteristic $0$ a polynomial is determined by its values on $\NN$,
by Corollary~\ref{rings:corollary-roots}, so the two descriptions agree.)
\end{proof}

\begin{example}[Binet's formula]\label{generating-functions:example-binet}
For the Fibonacci numbers, $\chi(t) = t^2 - t - 1 = (t - \varphi)(t - \bar{\varphi})$ over $\RR$, with $\varphi = \frac{1 + \sqrt 5}{2}$ (the golden ratio) and $\bar{\varphi} = \frac{1 - \sqrt 5}{2}$,
since $\varphi + \bar{\varphi} = 1$ and $\varphi\bar{\varphi} = -1$. So $F_n = \beta\varphi^n + \beta^{\prime}\bar{\varphi}^n$; the initial values $F_0 = 0$ and $F_1 = 1$ give $\beta^{\prime} = -\beta$ and
$\beta(\varphi - \bar{\varphi}) = \beta\sqrt 5 = 1$:
\[
F_n = \frac{\varphi^n - \bar{\varphi}^n}{\sqrt 5}, \qquad L_n = \varphi^n + \bar{\varphi}^n,
\]
the second in the same way from $L_0 = 2$, $L_1 = 1$. Since $|\bar{\varphi}| < 0.62$, $|\bar{\varphi}^n/\sqrt 5| < \frac12$ for all $n \geq 0$, so $F_n$ is the integer nearest to $\varphi^n/\sqrt 5$, and
$L_n$ is the integer nearest to $\varphi^n$ for $n \geq 2$. For instance $\varphi^{10}/\sqrt 5 = 55.0036\ldots$ and $F_{10} = 55$.
\end{example}

\begin{example}[A double root]\label{generating-functions:example-double-root}
The recurrence $a_{n+2} = 4a_{n+1} - 4a_n$ has $\chi = (t - 2)^2$, so its solutions are $a_n = (\beta + \gamma n)2^n$. With $a_0 = 1$ and $a_1 = 4$: $\beta = 1$, $2(1 + \gamma) = 4$, so
$a_n = (n + 1)2^n$, the coefficients of $\frac{1}{(1 - 2x)^2}$. The factor $n$ is the effect of the repeated root: the sequences $\beta2^n$ alone form a space of dimension $1$, while the solutions of a
recurrence of order $2$ form a space of dimension $2$.
\end{example}

\begin{counterexample}\label{generating-functions:counterexample-splitting}
The splitting hypothesis cannot be dropped over a field that is not algebraically closed. Over $\QQ$, the recurrence $a_{n+2} = -a_n$ has $\chi = t^2 + 1$, without rational roots, and the theorem says nothing about its solution
$0, 1, 0, -1, 0, 1, \ldots$ over $\QQ$; one must pass to the splitting field $\QQ(i)$, where $a_n = \frac{i^n - (-i)^n}{2i}$. And the hypothesis $c_d \neq 0$ excludes the root $0$: the recurrence $a_{n+2} = a_{n+1}$ has
the solutions $a_0$ arbitrary, $a_1 = a_2 = \cdots$, which are of the form $\beta \cdot 1^n$ only from $n = 1$ on; a root $0$ of multiplicity $m$ only affects the first $m$ terms.
\end{counterexample}

\subsection{Fibonacci identities}\label{generating-functions:subsection-fibonacci}

\begin{proposition}\label{generating-functions:proposition-fibonacci}
For $m, n \geq 0$ (and $n \geq 1$ where $F_{n-1}$ occurs):
\begin{enumerate}
\item $\begin{pmatrix}1 & 1\\ 1 & 0\end{pmatrix}^n = \begin{pmatrix}F_{n+1} & F_n\\ F_n & F_{n-1}\end{pmatrix}$ for $n \geq 1$;
\item (Cassini) $F_{n+1}F_{n-1} - F_n^2 = (-1)^n$;
\item (Addition) $F_{m+n} = F_{m+1}F_n + F_mF_{n-1}$, and $F_{2n} = F_nL_n$;
\item $F_m$ divides $F_{km}$ for every $k \geq 1$, and $\gcd(F_m, F_n) = F_{\gcd(m, n)}$;
\item $F_0 + F_1 + \cdots + F_n = F_{n+2} - 1$ and $F_1^2 + \cdots + F_n^2 = F_nF_{n+1}$.
\end{enumerate}
\end{proposition}

\begin{proof}
(1) By induction: for $n = 1$ it is the definition of the matrix ($F_0 = 0$, $F_1 = F_2 = 1$), and multiplying by the matrix on the right gives the rows $(F_{n+1} + F_n, F_{n+1})$ and
$(F_n + F_{n-1}, F_n)$, that is, $(F_{n+2}, F_{n+1})$ and $(F_{n+1}, F_n)$.
(2) Take determinants in (1): the determinant is multiplicative (Proposition~\ref{linear-algebra:proposition-determinant-properties}), and the matrix has determinant $-1$.
(3) The entry in the top right corner of $M^{m+n} = M^mM^n$, where $M$ is the matrix of (1), is $F_{m+1}F_n + F_mF_{n-1}$ (for $m, n \geq 1$; for $m = 0$ the identity reads $F_n = F_n$).
For $m = n$: $F_{2n} = F_n(F_{n+1} + F_{n-1})$, and $F_{n+1} + F_{n-1} = L_n$, since both sides satisfy the Fibonacci recurrence and agree for $n = 1, 2$ ($1 + 0 = 1$, $2 + 1 = 3$).
(4) By (3), $F_{(k+1)m} = F_{km+1}F_m + F_{km}F_{m-1}$, so $F_m \mid F_{km}$ by induction on $k$. Consecutive Fibonacci numbers are coprime, by (2). Now
$\gcd(F_{qn+r}, F_n) = \gcd(F_{qn+1}F_r + F_{qn}F_{r-1}, F_n) = \gcd(F_{qn+1}F_r, F_n) = \gcd(F_r, F_n)$ for $n \geq 1$ and $1 \leq r < n$: the second equality as $F_n \mid F_{qn}$, the third as
$\gcd(F_{qn+1}, F_n)$ divides $\gcd(F_{qn+1}, F_{qn}) = 1$ (for $q = 0$, $F_1 = 1$). So $\gcd(F_m, F_n) = \gcd(F_n, F_r)$ whenever $m = qn + r$ with $0 \leq r < n$ (for $r = 0$ both sides are
$F_n$, as $F_n \mid F_{qn}$ and $F_0 = 0$): the pair of indices follows the Euclidean algorithm, $\gcd(m, n) = \gcd(n, r)$, which ends at a pair $(d, 0)$ with $d = \gcd(m, n)$, and
$\gcd(F_d, F_0) = F_d$.
(5) Both are by induction: $F_{n+2} - 1 + F_{n+1} = F_{n+3} - 1$, and $F_nF_{n+1} + F_{n+1}^2 = F_{n+1}F_{n+2}$.
\end{proof}

\begin{remark}\label{generating-functions:remark-fibonacci-generating}
The sum $\sum_{k \leq n}F_k = F_{n+2} - 1$ is also immediate from generating functions: by Example~\ref{generating-functions:example-partial-sums} the partial sums have generating function
$\frac{x}{(1 - x)(1 - x - x^2)}$, and $\frac{x}{(1 - x)(1 - x - x^2)} = \frac{1}{x^2}\Bigl(\frac{x}{1 - x - x^2} - x\Bigr) - \frac{1}{1 - x}$, as one checks by clearing denominators; the
first term, a series of order $2$ divided by $x^2$, is $\sum_nF_{n+2}x^n$ by the shift rule of Lemma~\ref{generating-functions:lemma-ring}. Identity (4) shows that
$F_n$ can only be prime if $n$ is prime or $n = 4$ ($F_4 = 3$): if $n = ab$ with $3 \leq a < n$, then $1 < F_a < F_n$ and $F_a \mid F_n$, and the only composite $n$ without such a factor is $4$; the converse fails, $F_{19} = 4181 = 37 \cdot 113$, and it is unknown whether there are infinitely many Fibonacci primes.
\end{remark}

\subsection{Counting with automata: the transfer matrix}\label{generating-functions:subsection-transfer-matrix}

Many sequences satisfying linear recurrences count words, or walks, built step by step with a finite memory. The \emph{transfer matrix method} proves that such counting sequences are always
rational and gives the denominator.

\begin{theorem}[Transfer matrix]\label{generating-functions:theorem-transfer-matrix}
Let $A$ be a $k \times k$ matrix with entries in $R$; think of $A_{ij}$ as the number (or weight) of steps from state $i$ to state $j$. Then in the matrix ring $\mathrm{M}_k(R[[x]])$,
\[
\begin{aligned}
&\sum_{n \geq 0}A^nx^n = (I - xA)^{-1},\\
&\sum_{n \geq 0}(A^n)_{ij}x^n = \frac{(-1)^{i+j}\det\bigl((I - xA)_{\hat{\jmath}\hat{\imath}}\bigr)}{\det(I - xA)},
\end{aligned}
\]
where $(I - xA)_{\hat{\jmath}\hat{\imath}}$ is the matrix with row $j$ and column $i$ removed. In particular each generating function $\sum_n(A^n)_{ij}x^n$ is rational with denominator
$\det(I - xA) = 1 - \operatorname{tr}(A)x + \cdots + (-1)^k\det(A)x^k$, of degree at most $k$, and the sequence $((A^n)_{ij})_n$ satisfies a linear recurrence of order at most $k$.
\end{theorem}

\begin{proof}
The family $(A^nx^n)$ of matrices is summable entrywise, and $(I - xA)\sum_nA^nx^n = \sum_nA^nx^n - \sum_nA^{n+1}x^{n+1} = I$, which shows the first formula ($I - xA$ commutes with $\sum A^nx^n$, so the
inverse is two-sided). The determinant $\det(I - xA)$ is a polynomial in $x$ of degree at most $k$ with constant term $\det I = 1$, so it is a unit of $R[[x]]$, and the adjugate matrix satisfies
$M\operatorname{adj}(M) = \det(M)I$ over the commutative ring $R[[x]]$ (Proposition~\ref{linear-algebra:proposition-determinant-properties}(4)). So
$(I - xA)^{-1} = \det(I - xA)^{-1}\operatorname{adj}(I - xA)$, which is the second formula; its numerators are minors of size $k - 1$, polynomials of degree at most $k - 1$, and
Theorem~\ref{generating-functions:theorem-rational} gives the recurrence. The coefficient of $x$ in $\det(I - xA)$ is $-\operatorname{tr}(A)$ and that of $x^k$ is $\det(-A)$, by the Leibniz
expansion of the determinant.
\end{proof}

\begin{example}\label{generating-functions:example-no-consecutive-ones}
Let $w_n$ be the number of words of length $n$ in the letters $0, 1$ with no two consecutive letters $1$. Classify by the last letter: the state is $0$ (last letter $0$, or the empty word) or $1$. From
state $0$ one may append $0$ or $1$; from state $1$ only $0$. So $A = \begin{pmatrix}1 & 1\\ 1 & 0\end{pmatrix}$, $\det(I - xA) = 1 - x - x^2$, and $w_n$ is the sum of the entries of the first row of
$A^n$, which is $F_{n+1} + F_n = F_{n+2}$ by Proposition~\ref{generating-functions:proposition-fibonacci}(1): $w_0, w_1, \ldots = 1, 2, 3, 5, 8, 13, \ldots$. Similarly, the words of length $n$ in
the letters $a, b, c$ with no two consecutive letters $a$ have the transfer matrix $\begin{pmatrix}2 & 1\\ 2 & 0\end{pmatrix}$ (states: last letter not $a$, last letter $a$), and their number satisfies
$u_{n+2} = 2u_{n+1} + 2u_n$: $1, 3, 8, 22, 60, 164, \ldots$, with generating function $\frac{1 + x}{1 - 2x - 2x^2}$.
\end{example}

\begin{example}\label{generating-functions:example-closed-walks}
The number of closed walks of length $n$ in a finite graph is $\operatorname{tr}(A^n)$, where $A$ is the adjacency matrix. For the cycle with $3$ vertices, $A = J - I$ with $J$ the all-ones
matrix; $J^2 = 3J$, so $J$ has the eigenvalues $3, 0, 0$ and $A$ has $2, -1, -1$. In a Jordan basis (Theorem~\ref{linear-algebra:theorem-jordan}) $A^n$ is triangular with diagonal entries
$2^n, (-1)^n, (-1)^n$, so $\operatorname{tr}(A^n) = 2^n + 2(-1)^n$; and indeed $\det(I - xA) = (1 - 2x)(1 + x)^2$. Graph theory uses such counts to read properties of a graph
off the spectrum of its adjacency matrix.
\end{example}

\subsection{Making change}\label{generating-functions:subsection-change}

In how many ways can one pay $n$ cents with coins of $1$, $5$, $10$, $25$ and $50$ cents? This is P\'olya's problem; the answer is the coefficient of $x^n$ in
\[
\frac{1}{(1 - x)(1 - x^5)(1 - x^{10})(1 - x^{25})(1 - x^{50})},
\]
because $\frac{1}{1 - x^c} = \sum_mx^{cm}$ and a way of paying is a choice of the numbers $m_1, m_5, \ldots$ of coins of each kind with $m_1 + 5m_5 + 10m_{10} + 25m_{25} + 50m_{50} = n$. Multiplying the
factors one at a time is an efficient algorithm: if $a_n$ are the coefficients of a series $A$, the coefficients of $A/(1 - x^c)$ are $b_n = a_n + b_{n-c}$. It gives $292$ ways to pay a dollar
($n = 100$), and $242$ without the half-dollar.

More generally, for positive integers $c_1, \ldots, c_k$ the number of solutions of $c_1m_1 + \cdots + c_km_k = n$ in natural numbers has the rational generating function $\prod_i(1 - x^{c_i})^{-1}$,
whose denominator has the roots of unity of orders dividing the $c_i$ as reciprocal roots. By Theorem~\ref{generating-functions:theorem-closed-form} over $\CC$, the number is a \emph{quasi-polynomial}:
a sum of terms $p(n)\zeta^n$ with $\zeta$ a root of unity, that is, a function that agrees with a polynomial on each residue class modulo the least common multiple of the $c_i$. The leading term
comes from the root $1$, of multiplicity $k$, and is $\frac{n^{k-1}}{(k - 1)!\,c_1 \cdots c_k}$; it is the volume of a slice of a simplex, as the counting interpretation suggests.

\begin{example}\label{generating-functions:example-three-parts}
The number $a_n$ of partitions of $n$ into parts at most $3$ (equivalently, into at most $3$ parts, Exercise~\ref{combinatorics:exercise-partitions-three-parts}) is the integer nearest to
$(n + 3)^2/12$. Here is a proof by generating functions. Let $A = \frac{1}{(1 - x)(1 - x^2)(1 - x^3)}$. Since $1 - x^6 = (1 - x^3)(1 + x^3)$,
\[
\begin{aligned}
&(1 - x^6)A = \frac{1 + x^3}{(1 - x)(1 - x^2)},\\
&\frac{1}{(1 - x)(1 - x^2)} = \sum_n\Bigl(\Bigl\lfloor\frac{n}{2}\Bigr\rfloor + 1\Bigr)x^n,
\end{aligned}
\]
the latter because a solution of $m_1 + 2m_2 = n$ is determined by $m_2 \in \{0, \ldots, \lfloor n/2\rfloor\}$. So for $n \geq 6$, comparing coefficients of $x^n$, $a_n - a_{n-6} = \lfloor n/2\rfloor + 1 + \lfloor(n - 3)/2\rfloor + 1 = n$ (check $n = 2m$ and $n = 2m + 1$ separately). The function
$f(n) = (n + 3)^2/12$ satisfies $f(n) - f(n - 6) = \frac{(n + 3)^2 - (n - 3)^2}{12} = n$ as well. Hence $g(n) = a_n - f(n)$ satisfies $g(n) = g(n - 6)$ for $n \geq 6$: it is periodic with period $6$.
Its values for $n = 0, \ldots, 5$ (with $a_0, \ldots, a_5 = 1, 1, 2, 3, 4, 5$) are $\frac14, -\frac13, -\frac1{12}, 0, -\frac1{12}, -\frac13$, all of absolute value less than $\frac12$.
\end{example}

\section{Catalan numbers}\label{generating-functions:section-catalan}

The Catalan numbers $C_n = \frac{1}{n + 1}\binom{2n}{n}$, $1, 1, 2, 5, 14, 42, 132, \ldots$, were defined in Proposition~\ref{combinatorics:proposition-lattice-paths} as the numbers of lattice paths
that do not go above the diagonal. Stanley's book \cite{stanley-catalan} lists more than two hundred families of objects counted by them. They are the first example of a generating function defined
by a \emph{nonlinear} equation, and of the method that decomposes a structure into smaller ones of the same kind.

\subsection{Binary trees and the quadratic equation}\label{generating-functions:subsection-binary-trees}

\begin{definition}\label{generating-functions:definition-binary-tree}
The sets $\mathcal{B}_n$ of \emph{binary trees with $n$ nodes} are defined by recursion on $n$: $\mathcal{B}_0 = \{\varepsilon\}$ consists of the \emph{empty tree} $\varepsilon$, and for $n \geq 1$,
\[
\mathcal{B}_n = \bigsqcup_{i + j = n - 1}\mathcal{B}_i \times \mathcal{B}_j;
\]
a tree $(L, R) \in \mathcal{B}_i \times \mathcal{B}_j$ is drawn as a root node with the \emph{left subtree} $L$ and the \emph{right subtree} $R$ hanging below it. Let $b_n = |\mathcal{B}_n|$ and
$B = \sum_nb_nx^n$.
\end{definition}

The definition is a formal version of the picture: a binary tree is either empty, or a root with an ordered pair of binary trees attached. The order matters: the two trees with $2$ nodes, a root with a
left child and a root with a right child, are different.

\begin{theorem}\label{generating-functions:theorem-catalan}
$b_0 = 1$ and $b_n = \sum_{i=0}^{n-1}b_ib_{n-1-i}$ for $n \geq 1$ (Segner's recurrence). The generating function satisfies $B = 1 + xB^2$, and
\[
B = \frac{1 - \sqrt{1 - 4x}}{2x}, \qquad b_n = \frac{1}{n + 1}\binom{2n}{n} = C_n.
\]
\end{theorem}

\begin{proof}
\emph{The recurrence} is the rule of sum and the rule of product (Theorems~\ref{combinatorics:theorem-rule-of-sum} and~\ref{combinatorics:theorem-rule-of-product}) applied to the definition.
\emph{The equation.} The coefficient of $x^n$ in $1 + xB^2$ is $1$ for $n = 0$ and $\sum_{i+j=n-1}b_ib_j = b_n$ for $n \geq 1$.
\emph{Solving it.} Multiply $xB^2 - B + 1 = 0$ by $4x$ and complete the square: $(1 - 2xB)^2 = 1 - 4xB + 4x^2B^2 = 1 - 4x$. The series $1 - 2xB$ has constant term $1$ and its square is $1 - 4x$;
by Theorem~\ref{generating-functions:theorem-binomial-series}(4), $1 - 2xB = \sqrt{1 - 4x}$, where $\sqrt{1 - 4x} = (1 - 4x)^{1/2}$. (The other root of the quadratic, with $+\sqrt{1 - 4x}$, is excluded:
$1 + \sqrt{1 - 4x}$ has constant term $2$, so it is not $2xB$.) Hence $2xB = 1 - \sqrt{1 - 4x}$, a series without constant term, and by Example~\ref{generating-functions:example-sqrt},
$[x^{n+1}](1 - \sqrt{1 - 4x}) = \frac{2}{n + 1}\binom{2n}{n}$, so $b_n = \frac{1}{n + 1}\binom{2n}{n}$, which is $C_n$ by Proposition~\ref{combinatorics:proposition-lattice-paths}.
\end{proof}

The solution can be read in two ways. Directly, the closed formula for $b_n$ comes out of the binomial series for $(1 - 4x)^{1/2}$. Indirectly, Theorem~\ref{generating-functions:theorem-catalan}
shows that \emph{any} family of objects satisfying Segner's recurrence, with one object of size $0$, is counted by the Catalan numbers; it is enough to find a decomposition of the objects into an
ordered pair of smaller ones. We do this for several families, and denote the generating function by $C = \sum_nC_nx^n$.

\begin{corollary}\label{generating-functions:corollary-catalan-recurrence}
$C_{n+1} = \sum_{i=0}^{n}C_iC_{n-i}$ for $n \geq 0$; $C_{n+1} = \frac{2(2n + 1)}{n + 2}C_n$; and $C = \frac{1}{1 - xC}$, that is, $C$ is also the generating function of finite sequences of objects
counted by $xC$.
\end{corollary}

\begin{proof}
The first is Segner's recurrence. The second: $\frac{C_{n+1}}{C_n} = \frac{(2n + 2)!}{(n + 1)!(n + 2)!} \cdot \frac{n!(n + 1)!}{(2n)!} = \frac{(2n + 1)(2n + 2)}{(n + 1)(n + 2)}$. The third: $C = 1 + xC^2$
gives $C(1 - xC) = 1$, and $1 - xC$ is a unit.
\end{proof}

\subsection{Interpretations}\label{generating-functions:subsection-catalan-interpretations}

\begin{proposition}\label{generating-functions:proposition-catalan-interpretations}
Each of the following sets has $C_n$ elements.
\begin{enumerate}
\item \emph{Dyck words}: words with $n$ letters $U$ and $n$ letters $D$ such that every prefix has at least as many letters $U$ as letters $D$.
\item \emph{Bracketings} of a product $x_0x_1 \cdots x_n$ of $n + 1$ factors in a fixed order, by $n$ pairs of brackets, that is, the ways to compute it by $n$ multiplications of two factors.
\item \emph{Plane trees with $n$ edges}: here a plane tree is a root together with a finite (possibly empty) sequence of plane trees, its subtrees; the number of edges of a tree with subtrees
$T_1, \ldots, T_k$ is $k + \sum_i\mathrm{edges}(T_i)$.
\item \emph{Triangulations} of a convex polygon with $n + 2$ vertices, that is, maximal sets of pairwise non-crossing diagonals (Definition below).
\item \emph{Non-crossing partitions} of $[n]$: set partitions of $[n]$ in which there are no $a < b < c < d$ with $a, c$ in one block and $b, d$ in another.
\end{enumerate}
\end{proposition}

For (4) we use the following combinatorial model of a convex polygon, which avoids geometry. The vertices are $0, 1, \ldots, m - 1$, its \emph{sides} are the pairs $\{i, i + 1\}$ and $\{0, m - 1\}$,
and its \emph{diagonals} are the other pairs $\{i, j\}$ with $i < j$. Two pairs $\{a, b\}$ and $\{c, d\}$ with $a < b$, $c < d$ \emph{cross} if $a < c < b < d$ or $c < a < d < b$; this is what
chords of a circle through points in this cyclic order do. A \emph{triangulation} is a set of pairwise non-crossing diagonals, maximal for inclusion. For $m = 3$ the only triangulation is $\emptyset$.

\begin{proof}
In each case we show that the number $u_n$ of objects of size $n$ satisfies $u_0 = 1$ and Segner's recurrence, or (for (3) and (5)) that the generating function $U$ satisfies $U = 1/(1 - xU)$, which is
equivalent to $U = 1 + xU^2$ (Corollary~\ref{generating-functions:corollary-catalan-recurrence} shows one direction, and the other is the same computation backwards). In both cases $U = C$, because the
equation $U = 1 + xU^2$ determines the coefficients of $U$ recursively.

(1) The empty word is the only Dyck word of length $0$. A nonempty Dyck word $w$ starts with $U$; let $2k + 2$ be the length of its shortest nonempty prefix with as many $U$s as $D$s (its \emph{first
return}). Then $w = UvDv^{\prime}$, where $v$ and $v^{\prime}$ are Dyck words of lengths $2k$ and $2(n - 1 - k)$: every prefix of $v$ has at least as many $U$s as $D$s, because a prefix $Uv_0$ of
$w$ with $v_0$ a prefix of $v$ has \emph{more} $U$s than $D$s (by minimality of the first return), and $v^{\prime}$ starts at a point where the counts are equal. Conversely $UvDv^{\prime}$ is a Dyck word
whose first return is after $UvD$. This bijection with the pairs $(v, v^{\prime})$ gives Segner's recurrence. (This also reproves the count of
Proposition~\ref{combinatorics:proposition-lattice-paths}, with $U = E$ and $D = N$.)

(2) A product of one factor has one bracketing. For $n \geq 1$, the last multiplication computes $(x_0 \cdots x_i)(x_{i+1} \cdots x_n)$ for a unique $i \in \{0, \ldots, n - 1\}$, and the two factors are
bracketed independently, with $i$ and $n - 1 - i$ multiplications.

(3) The plane trees with $n$ edges whose root has $k$ subtrees $T_1, \ldots, T_k$ correspond to sequences of $k$ plane trees with $n - k$ edges in all. So the generating function $P$ (by edges) is
$P = \sum_{k \geq 0}(xP)^k = \frac{1}{1 - xP}$, the family $((xP)^k)_k$ being summable since $\operatorname{ord}((xP)^k) \geq k$.

(4) Let $m = n + 2$ and let $T$ be a triangulation; put $S = T \cup \{\text{sides}\}$. Let $k$ be the largest vertex $\leq n$ with $\{0, k\} \in S$ (it exists, since $\{0, 1\}$ is a side).
\emph{Claim: $\{k, n + 1\} \in S$.} If $k = n$ it is a side. Otherwise $\{k, n + 1\}$ is a diagonal, and it crosses no element $\{a, b\}$ ($a < b$) of $S$: a crossing would need $a < k < b < n + 1$
(the other pattern needs $b > n + 1$). If $a = 0$, then $\{0, b\} \in S$ with $k < b \leq n$, contradicting the choice of $k$. If $a > 0$, then $k \geq 2$ and $\{0, k\}$ is a diagonal of $T$ crossed by
$\{a, b\}$, since $0 < a < k < b$; but $\{a, b\}$ is a side or a diagonal of $T$, and sides cross nothing ($a < c < b < d$ is impossible when $b = a + 1$, and $\{0, n + 1\}$ crosses nothing since no vertex
lies outside $[0, n + 1]$). By maximality, $\{k, n + 1\} \in T$.
\emph{Splitting.} So the triangle $\{0, k, n + 1\}$ has its three sides in $S$; and every other diagonal of $T$ has both ends in $\{0, \ldots, k\}$ or both in $\{k, \ldots, n + 1\}$, since otherwise it
would cross $\{0, k\}$ or $\{k, n + 1\}$. Say that a polygon with $s + 2$ vertices has \emph{size} $s$, and let the ``polygon'' with $2$ vertices (size $0$) have one, empty, triangulation. The diagonals of $T$ other than $\{0, k\}$ and
$\{k, n + 1\}$ form triangulations of the polygons with vertices $0, \ldots, k$ (size $k - 1$) and $k, \ldots, n + 1$ (size $n - k$), relabelled in order. Maximality in a small polygon follows from
maximality of $T$: a diagonal of the small polygon that crosses nothing there crosses nothing in $T$, since the diagonals of the other small polygon, $\{0, k\}$ and $\{k, n + 1\}$ have both ends
outside its open range of vertices. Conversely, a choice of $k \in \{1, \ldots, n\}$ and of triangulations of the two smaller polygons, together with the diagonals $\{0, k\}$ (if $k \geq 2$) and
$\{k, n + 1\}$ (if $k \leq n - 1$), gives a triangulation: a further diagonal $\{i, j\}$ either lies in one small polygon, where it crosses something by maximality there, or has $i < k < j$, and then
it crosses $\{0, k\}$ (if $i > 0$) or $\{k, n + 1\}$ (if $i = 0$, so $j \leq n$). The maximal $k$ for this triangulation is the chosen one, since no diagonal $\{0, j\}$ with $j > k$ was added. This gives $u_n = \sum_{k=1}^{n}u_{k-1}u_{n-k}$, Segner's recurrence.

(5) Let $N$ be the generating function of the non-crossing partitions, with $N_0 = 1$. In a non-crossing partition of $[n]$, $n \geq 1$, let the block of $1$ be $\{1 = e_1 < e_2 < \cdots < e_r\}$, and
put $e_{r+1} = n + 1$. The \emph{gaps} $G_i = \{e_i + 1, \ldots, e_{i+1} - 1\}$, $1 \leq i \leq r$, partition $[n] \setminus \{e_1, \ldots, e_r\}$, and every other block lies inside one gap: a block
meeting two gaps $G_i$ and $G_j$, $i < j$, in elements $b \in G_i$ and $d \in G_j$ would cross the block of $1$, since $e_i < b < e_{i+1} \leq e_j < d$, which is the forbidden pattern with
$a = e_i$ and $c = e_{i+1}$ (here $i + 1 \leq j \leq r$, so $e_{i+1}$ is an element of the block of $1$). So the partition is the block of $1$ together with non-crossing partitions of the $r$ gaps (relabelled in order), and conversely every such choice is non-crossing,
since blocks in different gaps, or a block in a gap and the block of $1$, do not interleave. The gaps have sizes $s_1, \ldots, s_r \geq 0$ with $r + s_1 + \cdots + s_r = n$. Therefore
$N = \sum_{r \geq 0}(xN)^r = \frac{1}{1 - xN}$.
\end{proof}

\begin{remark}\label{generating-functions:remark-catalan-bijections}
The proof gives explicit recursive bijections between any two of these families, since they are all built by the same decomposition. For instance a binary tree $(L, R)$ corresponds to the Dyck
word $U\,w(L)\,D\,w(R)$, and to the bracketing $(b(L))(b(R))$ with $b(\varepsilon) = x$. Non-recursive descriptions are also known: the plane tree with $n$ edges corresponds to the Dyck word
obtained by walking around it from the root, writing $U$ for each edge traversed downwards and $D$ upwards.
\end{remark}

\begin{example}\label{generating-functions:example-catalan-small}
For $n = 3$ there are $C_3 = 5$ objects of each kind: the Dyck words $UUUDDD$, $UUDUDD$, $UUDDUD$, $UDUUDD$, $UDUDUD$; the bracketings $((x_0x_1)x_2)x_3$, $(x_0(x_1x_2))x_3$, $(x_0x_1)(x_2x_3)$,
$x_0((x_1x_2)x_3)$, $x_0(x_1(x_2x_3))$; the $5$ triangulations of a pentagon, each consisting of the two diagonals from one vertex; and all $5$ partitions of $[3]$, since a crossing needs $4$
elements. For $n = 4$, of the $15$ partitions of $[4]$ (the Bell number $B_4$) exactly one, $\{\{1, 3\}, \{2, 4\}\}$, is crossing, and indeed $C_4 = 14$.
\end{example}

\begin{counterexample}\label{generating-functions:counterexample-catalan-order}
Order matters in the definitions. If the two subtrees of a binary tree are not ordered (rooted trees in which each vertex has $0$ or $2$ unordered children), the numbers of such trees with
$1, 3, 5, 7, 9$ vertices are $1, 1, 1, 2, 3$ (the Wedderburn--Etherington numbers), not Catalan numbers: the decomposition is into an \emph{unordered} pair, and the rule of product no longer applies
directly. Likewise, the number of all set partitions of $[n]$ is the Bell number $B_n$ (Proposition~\ref{combinatorics:proposition-bell}), which exceeds $C_n$ from $n = 4$ on: the non-crossing condition
is what makes the gaps independent.
\end{counterexample}

\subsection{Parity and growth of the Catalan numbers}\label{generating-functions:subsection-catalan-arithmetic}

\begin{proposition}\label{generating-functions:proposition-catalan-parity}
$C_n$ is odd if and only if $n = 2^a - 1$ for some $a \geq 0$.
\end{proposition}

\begin{proof}
By Kummer's theorem (Theorem~\ref{combinatorics:theorem-kummer}) with $p = 2$, $v_2\binom{2n}{n} = 2s_2(n) - s_2(2n) = s_2(n)$, the number of digits $1$ in the binary expansion of $n$, since doubling
only shifts the digits. Write $n + 1 = 2^ab$ with $b$ odd; then $v_2(n + 1) = a$, and $n = 2^ab - 1$ ends with exactly $a$ binary digits $1$ (followed, to the left, by the digits of $b - 1$, which is even).
So $s_2(n) = a + s_2(b - 1) \geq a$, with equality if and only if $b = 1$. Since $v_2(C_n) = v_2\binom{2n}{n} - v_2(n + 1) = s_2(n) - a$, the number $C_n$ is odd exactly when $n = 2^a - 1$.
\end{proof}

So $C_0, C_1, C_3, C_7 = 1, 1, 5, 429$ are odd, and $C_2 = 2$, $C_4 = 14$, $C_5 = 42$, $C_6 = 132$ are even. The growth of $C_n$ is that of $\binom{2n}{n}$ divided by $n + 1$, roughly $4^n/n^{3/2}$;
Section~\ref{generating-functions:section-estimates} proves the elementary bounds.

\section{Partitions and infinite products}\label{generating-functions:section-partitions}

Partitions of integers (Section~\ref{combinatorics:section-partitions}) are where generating functions first showed their power, in the work of Euler. Their generating functions are infinite
products, and identities between infinite products prove theorems about partitions that are hard to see by counting.

\subsection{Rules for infinite products}\label{generating-functions:subsection-product-rules}

\begin{lemma}\label{generating-functions:lemma-product-rules}
Let $(g_k)_{k \geq 1}$ and $(h_k)_{k \geq 1}$ be families satisfying the hypotheses of Definition~\ref{generating-functions:definition-infinite-product}, and write $\prod_k$ for $\prod_{k \geq 1}$.
\begin{enumerate}
\item (Factorwise products) $\prod_k(1 + g_k) \cdot \prod_k(1 + h_k) = \prod_k(1 + g_k)(1 + h_k)$, where the right-hand side is the infinite product of the family $(g_k + h_k + g_kh_k)_k$.
\item (Inverses) Each $1 + g_k$ is a unit, the family $((1 + g_k)^{-1} - 1)_k$ satisfies the hypotheses, and $\prod_k(1 + g_k)^{-1} = \bigl(\prod_k(1 + g_k)\bigr)^{-1}$.
\item (Reindexing and splitting) If $\pi$ is a bijection of $\{1, 2, \ldots\}$, then $\prod_k(1 + g_{\pi(k)}) = \prod_k(1 + g_k)$; and if $\{1, 2, \ldots\} = I \sqcup J$, then
$\prod_k(1 + g_k) = \prod_{k \in I}(1 + g_k)\prod_{k \in J}(1 + g_k)$.
\item (Logarithmic derivative) $D\prod_k(1 + g_k) = \prod_k(1 + g_k) \cdot \sum_k\frac{Dg_k}{1 + g_k}$, the sum being summable.
\item (Coefficient homomorphisms) If $\phi \colon R \to S$ is a ring homomorphism and $\phi_\ast \colon R[[x]] \to S[[x]]$ applies $\phi$ to every coefficient, then $\phi_\ast$ is a ring homomorphism
that preserves summable sums, limits and infinite products. So does the substitution $f \mapsto f(x^m)$, $m \geq 1$.
\end{enumerate}
\end{lemma}

\begin{proof}
All the statements follow from three facts: infinite products are limits of their finite partial products (Proposition~\ref{generating-functions:proposition-infinite-product}(3)); finite products
can be rearranged freely; and limits respect sums and products (the remark after Definition~\ref{generating-functions:definition-convergence}), as well as the operations $D$, $\phi_\ast$ and
$f \mapsto f(x^m)$, because each coefficient of the result depends on finitely many coefficients of the argument, and $\operatorname{ord}(f(x^m)) = m\operatorname{ord}(f)$.
(1) The family $(g_k + h_k + g_kh_k)$ satisfies the hypotheses, since $\operatorname{ord}(g_k + h_k + g_kh_k) \geq \min(\operatorname{ord} g_k, \operatorname{ord} h_k)$. The partial products of
the right-hand side are the products of the partial products of the two factors on the left; pass to the limit.
(2) $1 + g_k$ is a unit since its constant term is $1$, and $(1 + g_k)^{-1} - 1 = -g_k(1 + g_k)^{-1}$ has the same order as $g_k$. By (1), $\prod_k(1 + g_k) \cdot \prod_k(1 + g_k)^{-1} = \prod_k1 = 1$.
(3) The partial products over $\{1, \ldots, N\}$ and over $\pi^{-1}(\{1, \ldots, N\})$ eventually agree with the infinite product in any given coefficient, by the proof of
Proposition~\ref{generating-functions:proposition-infinite-product}(3), which works for any increasing sequence of finite sets exhausting the index set. The splitting is
Proposition~\ref{generating-functions:proposition-infinite-product}(4).
(4) For a finite product this is the Leibniz rule, $D\prod_{k \leq N}(1 + g_k) = \prod_{k \leq N}(1 + g_k)\sum_{k \leq N}Dg_k/(1 + g_k)$. The family $(Dg_k/(1 + g_k))_k$ is summable, since
$\operatorname{ord}(Dg_k) \geq \operatorname{ord}(g_k) - 1$. Pass to the limit $N \to \infty$.
(5) $\phi_\ast$ is additive, and multiplicative because $\phi(\sum_{i+j=n}a_ib_j) = \sum_{i+j=n}\phi(a_i)\phi(b_j)$; it does not increase orders, and acts coefficient by coefficient, so it preserves
summable sums and limits, hence infinite products. The substitution $f \mapsto f(x^m)$ is the composition with $x^m$ (Lemma~\ref{generating-functions:lemma-composition}).
\end{proof}

\subsection{The partition generating function}\label{generating-functions:subsection-partition-gf}

\begin{theorem}[Euler]\label{generating-functions:theorem-partition-gf}
\[
\sum_{n \geq 0}p(n)x^n = \prod_{k \geq 1}\frac{1}{1 - x^k}.
\]
More generally, for a set $S$ of positive integers, $\prod_{k \in S}\frac{1}{1 - x^k}$ is the generating function of the partitions with all parts in $S$, and $\prod_{k \in S}(1 + x^k)$ that of the
partitions into distinct parts in $S$; and in $\ZZ[y][[x]]$,
\[
\prod_{k \geq 1}\frac{1}{1 - yx^k} = \sum_{n, m \geq 0}p_m(n)\,y^mx^n,
\]
where $p_m(n)$ is the number of partitions of $n$ into exactly $m$ parts.
\end{theorem}

\begin{proof}
A partition is determined by its \emph{multiplicities}: the number $m_k$ of parts equal to $k$, for $k \geq 1$, which is a finitely supported family of natural numbers, and the partition has size
$\sum_kkm_k$ and $\sum_km_k$ parts; conversely every finitely supported family $(m_k)$ comes from exactly one partition. Write $\frac{1}{1 - yx^k} = 1 + g_k$ with $g_k = \sum_{m \geq 1}y^mx^{km}$,
of order $k$ in $x$. By Definition~\ref{generating-functions:definition-infinite-product} and Proposition~\ref{generating-functions:proposition-summable}(3), the product is
\[
\sum_{S \text{ finite}}\prod_{k \in S}\sum_{m_k \geq 1}y^{m_k}x^{km_k} = \sum_{(m_k)}y^{\sum_km_k}x^{\sum_kkm_k},
\]
the sum over all finitely supported families $(m_k)_{k \geq 1}$ (the set $S$ is their support). The coefficient of $y^mx^n$ counts the partitions of $n$ with $m$ parts. Setting $y = 1$
(Lemma~\ref{generating-functions:lemma-product-rules}(5), with $\phi \colon \ZZ[y] \to \ZZ$, $y \mapsto 1$) gives the first formula. Restricting the factors to $k \in S$ restricts the parts to $S$,
and for distinct parts each $m_k$ is $0$ or $1$, which gives the factors $1 + x^k$.
\end{proof}

\begin{example}\label{generating-functions:example-partition-gf-variants}
\begin{enumerate}
\item Partitions into parts at most $m$: $\prod_{k=1}^{m}(1 - x^k)^{-1}$. By conjugation (Proposition~\ref{combinatorics:proposition-conjugate}), this also counts the partitions into at most $m$
parts.
\item Partitions in which no part occurs more than $d - 1$ times: $\prod_k(1 + x^k + \cdots + x^{(d-1)k}) = \prod_k\frac{1 - x^{dk}}{1 - x^k}$.
\item Partitions into odd parts: $\prod_{k \text{ odd}}(1 - x^k)^{-1}$. Partitions into distinct parts: $\prod_k(1 + x^k)$ (Example~\ref{generating-functions:example-product-distinct}).
\end{enumerate}
\end{example}

\begin{corollary}[Euler, Glaisher]\label{generating-functions:corollary-glaisher}
For $d \geq 2$, the number of partitions of $n$ with no part divisible by $d$ equals the number of partitions of $n$ in which no part occurs $d$ or more times. For $d = 2$: partitions into odd parts
and partitions into distinct parts are equinumerous.
\end{corollary}

\begin{proof}
By Example~\ref{generating-functions:example-partition-gf-variants}(2) and Lemma~\ref{generating-functions:lemma-product-rules}, the second number has the generating function
\[
\prod_{k \geq 1}\frac{1 - x^{dk}}{1 - x^k} = \frac{\prod_{k \geq 1}(1 - x^{dk})}{\prod_{d \mid k}(1 - x^k)\prod_{d \nmid k}(1 - x^k)} = \prod_{d \nmid k}\frac{1}{1 - x^k},
\]
using (1) and (2) of the lemma to write the product of quotients as a quotient of products, (3) to split $\prod_k(1 - x^k)$ according to whether $d \mid k$, and reindexing to identify
$\prod_{k \geq 1}(1 - x^{dk})$ with $\prod_{d \mid k}(1 - x^k)$. The last product is the generating function of partitions without parts divisible by $d$, by Theorem~\ref{generating-functions:theorem-partition-gf}.
\end{proof}

For $d = 2$ this is Theorem~\ref{combinatorics:theorem-odd-distinct}, proved there by Glaisher's bijection; here it is a two-line computation, and the bijection can be recovered from it, the factor
$\frac{1 - x^{2k}}{1 - x^k} = 1 + x^k$ corresponding to the binary expansion of multiplicities. This is the typical relation between the two methods: generating functions find and prove identities,
bijections explain them.

\subsection{Gaussian binomial coefficients}\label{generating-functions:subsection-gaussian}

\begin{definition}\label{generating-functions:definition-gaussian}
For $0 \leq k \leq n$, the \emph{Gaussian binomial coefficient} (or \emph{$q$-binomial coefficient}) is the polynomial
\[
\binom{n}{k}_q = \sum_{\lambda}q^{|\lambda|} \in \ZZ[q],
\]
the sum over the partitions $\lambda$ with at most $k$ parts, all at most $n - k$, that is, whose Ferrers diagram fits in a box with $k$ rows and $n - k$ columns; $|\lambda|$ is the sum of the parts.
We put $\binom{n}{k}_q = 0$ for $k < 0$ or $k > n$. For $m \geq 0$ we write $(q)_m = (1 - q)(1 - q^2) \cdots (1 - q^m)$, with $(q)_0 = 1$.
\end{definition}

At $q = 1$ the definition counts the partitions in the box, which correspond to the lattice paths from one corner of the box to the other (the boundary of the diagram), so
$\binom{n}{k}_1 = \binom{n}{k}$ by Proposition~\ref{combinatorics:proposition-lattice-paths}. The Gaussian coefficient refines the binomial coefficient by the area under the path.

\begin{proposition}\label{generating-functions:proposition-gaussian}
For $1 \leq k \leq n$:
\begin{enumerate}
\item ($q$-Pascal) $\binom{n}{k}_q = \binom{n - 1}{k - 1}_q + q^k\binom{n - 1}{k}_q = \binom{n - 1}{k}_q + q^{n-k}\binom{n - 1}{k - 1}_q$;
\item $\binom{n}{k}_q = \binom{n}{n - k}_q$, the polynomial has degree $k(n - k)$, and its coefficients are symmetric: $q^{k(n-k)}\binom{n}{k}_{1/q} = \binom{n}{k}_q$;
\item (product formula) in $\ZZ[[q]]$, and hence in $\ZZ[q]$ after clearing denominators,
\[
\binom{n}{k}_q = \frac{(q)_n}{(q)_k(q)_{n-k}} = \frac{(1 - q^n)(1 - q^{n-1}) \cdots (1 - q^{n-k+1})}{(1 - q)(1 - q^2) \cdots (1 - q^k)}.
\]
\end{enumerate}
\end{proposition}

\begin{proof}
(1) Let $\lambda$ fit in the $k \times (n - k)$ box. If $\lambda$ has fewer than $k$ parts, it fits in the $(k - 1) \times (n - k)$ box, and these $\lambda$ give $\binom{n - 1}{k - 1}_q$. If it has
exactly $k$ parts, subtracting $1$ from each part gives a partition in the $k \times (n - k - 1)$ box, of size $|\lambda| - k$, bijectively; these give $q^k\binom{n - 1}{k}_q$. (For $k = n$ the second
case is empty, consistently with $\binom{n-1}{n}_q = 0$.) The second identity is proved in the same way, distinguishing whether the largest part is less than $n - k$ or equal to it (and then removing
it).
(2) Conjugation (Proposition~\ref{combinatorics:proposition-conjugate}) exchanges the $k \times (n - k)$ and $(n - k) \times k$ boxes preserving sizes. The largest size is that of the full box,
$k(n - k)$, attained once. Taking the complement of a diagram in the box (rotated by a half turn) is a bijection between partitions of size $s$ and of size $k(n - k) - s$, which gives the symmetry.
(3) Subtracting the two forms of (1): $(1 - q^{n-k})\binom{n - 1}{k - 1}_q = (1 - q^k)\binom{n - 1}{k}_q$, so with $m = n - 1$, $\binom{m}{k}_q = \frac{1 - q^{m-k+1}}{1 - q^k}\binom{m}{k - 1}_q$ for
$1 \leq k \leq m$ (both sides vanish for $k = m + 1$, where the factor $1 - q^0 = 0$ appears). By induction on $k$ from $\binom{m}{0}_q = 1$ this gives the product formula, the inverses existing in
$\ZZ[[q]]$ since the constant terms are $1$. Multiplying out, $\binom{n}{k}_q(q)_k(q)_{n-k} = (q)_n$ is an identity of polynomials.
\end{proof}

\begin{example}\label{generating-functions:example-gaussian}
$\binom{4}{2}_q = \frac{(1 - q^4)(1 - q^3)}{(1 - q)(1 - q^2)} = (1 + q^2)(1 + q + q^2) = 1 + q + 2q^2 + q^3 + q^4$: the partitions in a $2 \times 2$ box are $\emptyset$, $1$, $2$, $1 + 1$, $2 + 1$,
$2 + 2$, of sizes $0, 1, 2, 2, 3, 4$. At $q = 1$ the value is $6 = \binom{4}{2}$.
\end{example}

\begin{theorem}\label{generating-functions:theorem-subspaces}
Let $q$ be a prime power and $\mathbb{F}_q$ the field with $q$ elements (Theorem~\ref{fields:theorem-finite-fields}). The number of $k$-dimensional subspaces of $\mathbb{F}_q^n$ is the value of
the polynomial $\binom{n}{k}_q$ at $q$.
\end{theorem}

\begin{proof}
Count the sequences $(v_1, \ldots, v_k)$ of linearly independent vectors in $\mathbb{F}_q^n$: $v_1$ can be any of the $q^n - 1$ nonzero vectors, and once $v_1, \ldots, v_{i}$ are chosen, $v_{i+1}$
can be any vector outside their span, which has $q^i$ elements: $q^n - q^i$ choices. So there are $\prod_{i=0}^{k-1}(q^n - q^i)$ such sequences (Corollary~\ref{combinatorics:corollary-successive-choices}).
Each $k$-dimensional subspace $W$ arises from exactly $\prod_{i=0}^{k-1}(q^k - q^i)$ of them, its ordered bases (the same count inside $W \cong \mathbb{F}_q^k$). The number of subspaces is the
quotient
\[
\frac{\prod_{i=0}^{k-1}q^i(q^{n-i} - 1)}{\prod_{i=0}^{k-1}q^i(q^{k-i} - 1)} = \frac{(q^n - 1)(q^{n-1} - 1) \cdots (q^{n-k+1} - 1)}{(q^k - 1) \cdots (q - 1)},
\]
which is the value at $q$ of the polynomial $\binom{n}{k}_q$, by the polynomial identity of Proposition~\ref{generating-functions:proposition-gaussian}(3) (both numerator and denominator change sign
$k$ times when written with factors $1 - q^j$).
\end{proof}

\begin{remark}\label{generating-functions:remark-schubert}
There is also a bijective proof, which explains the partitions: every $k$-dimensional subspace is the row space of a unique $k \times n$ matrix in reduced row echelon form
(Theorem~\ref{linear-algebra:theorem-gaussian-elimination}). If the pivots are in columns $j_1 < \cdots < j_k$, the free entries are those to the right of each pivot in its row and not in a pivot
column; the number of free entries in row $i$ is $n - k - (j_i - i)$, and the numbers $\lambda_i = n - k - (j_i - i)$ form a partition in the $k \times (n - k)$ box. Each such matrix form gives
$q^{|\lambda|}$ subspaces. In algebraic geometry these sets of subspaces are the Schubert cells of the Grassmannian.
\end{remark}

\begin{theorem}[$q$-binomial theorem]\label{generating-functions:theorem-q-binomial}
In $\ZZ[q, z]$, for $n \geq 0$,
\[
\prod_{i=0}^{n-1}(1 + q^iz) = \sum_{k=0}^{n}q^{k(k-1)/2}\binom{n}{k}_qz^k.
\]
\end{theorem}

\begin{proof}
By induction on $n$; for $n = 0$ both sides are $1$. Multiplying the identity for $n - 1$ by $1 + q^{n-1}z$, the coefficient of $z^k$ on the left becomes
$q^{k(k-1)/2}\binom{n - 1}{k}_q + q^{n-1}q^{(k-1)(k-2)/2}\binom{n - 1}{k - 1}_q$. Since $n - 1 + \frac{(k-1)(k-2)}{2} = \frac{k(k-1)}{2} + (n - k)$, this is
$q^{k(k-1)/2}\bigl(\binom{n - 1}{k}_q + q^{n-k}\binom{n - 1}{k - 1}_q\bigr) = q^{k(k-1)/2}\binom{n}{k}_q$ by the second form of $q$-Pascal.
\end{proof}

At $q = 1$ this is the binomial theorem; combinatorially, the coefficient of $z^kq^m$ on the left counts the $k$-element subsets of $\{0, \ldots, n - 1\}$ with sum $m$, and subtracting
$0, 1, \ldots, k - 1$ from the elements, in increasing order, turns such a subset into a partition in the $k \times (n - k)$ box, of size $m - k(k-1)/2$.

\begin{corollary}[Euler's identity]\label{generating-functions:corollary-euler-identities}
In $\ZZ[z][[q]]$,
\[
\prod_{i \geq 0}(1 + q^iz) = \sum_{k \geq 0}\frac{q^{k(k-1)/2}z^k}{(q)_k},
\]
where the left-hand side means $(1 + z)\prod_{i \geq 1}(1 + q^iz)$, the infinite product being that of Definition~\ref{generating-functions:definition-infinite-product} in $\ZZ[z][[q]]$.
\end{corollary}

\begin{proof}
For fixed $k$, $\binom{n}{k}_q \to \frac{1}{(q)_k}$ as $n \to \infty$: the coefficient of $q^m$ in $\binom{n}{k}_q$ is the number of partitions of $m$ into at most $k$ parts, all at most $n - k$,
which for $n - k \geq m$ is the number of partitions of $m$ into at most $k$ parts, the coefficient of $q^m$ in $\frac{1}{(q)_k} = \prod_{j=1}^{k}\frac{1}{1 - q^j}$
(Example~\ref{generating-functions:example-partition-gf-variants}(1)). In the sum of Theorem~\ref{generating-functions:theorem-q-binomial}, the coefficient of $q^m$ only involves the $k$ with
$k(k-1)/2 \leq m$, finitely many; so the right-hand side converges to $\sum_kq^{k(k-1)/2}z^k/(q)_k$, while the left-hand side, $(1 + z)$ times a partial product, converges to the infinite product.
\end{proof}

\begin{counterexample}\label{generating-functions:counterexample-two-variables}
The companion identity $\prod_{i \geq 0}(1 - q^iz)^{-1} = \sum_{k \geq 0}z^k/(q)_k$ does not make sense in $\ZZ[z][[q]]$ as it stands: the factor $(1 - z)^{-1}$ is not a polynomial in $z$, and on the
right the coefficient of $q^0$ receives the contributions $z^k$ from every $k$, so the family is not summable in the sense of Definition~\ref{generating-functions:definition-summable}. It holds in a
ring of series in two variables with a summability condition on monomials $z^kq^m$, which we do not need. The choice of the ring is part of the statement of a formal identity.
\end{counterexample}

\subsection{The pentagonal number theorem}\label{generating-functions:subsection-pentagonal}

Expanding $\prod_{k \geq 1}(1 - x^k)$, Euler found the astonishingly sparse series $1 - x - x^2 + x^5 + x^7 - x^{12} - x^{15} + x^{22} + x^{26} - \cdots$. The exponents are the
\emph{pentagonal numbers} $\omega(j) = j(3j - 1)/2$ for $j \in \ZZ$: $0, 1, 2, 5, 7, 12, 15, 22, 26, \ldots$ for $j = 0, 1, -1, 2, -2, 3, -3, \ldots$.

\begin{theorem}[Euler's pentagonal number theorem]\label{generating-functions:theorem-pentagonal}
\[
\prod_{k \geq 1}(1 - x^k) = \sum_{j \in \ZZ}(-1)^jx^{j(3j-1)/2}.
\]
\end{theorem}

\begin{proof}
\emph{Step 1: signed count.} By Definition~\ref{generating-functions:definition-infinite-product} with $g_k = -x^k$, the left-hand side is $\sum_S(-1)^{|S|}x^{\sum S}$ over finite sets $S$ of positive
integers, so its coefficient of $x^n$ is $e(n) - o(n)$, where $e(n)$ and $o(n)$ are the numbers of partitions of $n$ into an even, resp.\ odd, number of distinct parts.

\emph{Step 2: the pentagonal numbers are distinct.} For $k \geq 1$, $\omega(k) = \frac{k(3k-1)}{2} < \omega(-k) = \frac{k(3k+1)}{2} < \omega(k + 1) = \frac{(k+1)(3k+2)}{2}$, so $\omega$ is injective
on $\ZZ$, and it suffices to show: $e(n) - o(n) = (-1)^j$ if $n = \omega(j)$, and $0$ if $n$ is not pentagonal.

\emph{Step 3: Franklin's involution.} Let $\lambda = (\lambda_1 > \cdots > \lambda_k)$ be a partition into $k \geq 1$ distinct parts. Its \emph{base} is its smallest part $s = \lambda_k$, and its
\emph{slope} is the largest $\sigma$ such that $\lambda_1, \lambda_2, \ldots, \lambda_\sigma$ are consecutive integers ($\lambda_i = \lambda_1 - i + 1$ for $i \leq \sigma$); so $1 \leq \sigma \leq k$. In
the Ferrers diagram, the base is the bottom row and the slope is the longest diagonal of dots running down at $45$ degrees from the end of the top row. Call $\lambda$ \emph{exceptional} if $\sigma = k$
and $s \in \{k, k + 1\}$, that is, if the slope and the base share a dot and the move below cannot be performed. Define $\phi(\lambda)$ for non-exceptional $\lambda$:
\begin{itemize}
\item \emph{Case A, $s \leq \sigma$}: remove the base and add $1$ to each of the $s$ largest parts. (Since $\lambda$ is not exceptional, $s < k$, so there are at least $s$ other parts.)
\item \emph{Case B, $s > \sigma$}: subtract $1$ from each of the $\sigma$ largest parts and add a new smallest part $\sigma$.
\end{itemize}

\emph{Step 4: $\phi$ is well defined and exchanges the cases.} In case A the result is $\mu = (\lambda_1 + 1, \ldots, \lambda_s + 1, \lambda_{s+1}, \ldots, \lambda_{k-1})$, with $k - 1$ parts. Its
parts are distinct, as $\lambda_s + 1 > \lambda_{s+1} + 1 > \lambda_{s+1}$. Its slope is exactly $s$: the first $s$ parts are consecutive (as $s \leq \sigma$), and if $s + 1 \leq k - 1$, then
$\mu_{s+1} = \lambda_{s+1} \leq \lambda_s - 1 = \mu_s - 2$. Its base is $\lambda_{k-1}$ or $\lambda_{k-1} + 1$, in both cases greater than $\lambda_k = s$. So $\mu$ falls in case B. It is not
exceptional: an exceptional partition in case B has as many parts as its slope and base equal to slope plus one; here that would mean $s = k - 1$ and base $k$, but then the base of $\mu$ is
$\lambda_{k-1} + 1 \geq \lambda_k + 2 = k + 1$.

In case B the result is $\mu = (\lambda_1 - 1, \ldots, \lambda_\sigma - 1, \lambda_{\sigma+1}, \ldots, \lambda_k, \sigma)$, with $k + 1$ parts. Its parts are distinct and positive: if $\sigma < k$,
then $\lambda_{\sigma+1} \leq \lambda_\sigma - 2$ by the maximality of the slope, and $\sigma < s = \lambda_k$; if $\sigma = k$, non-exceptionality gives $s \geq k + 2$, so the parts $\lambda_i - 1$ are
all at least $s - 1 > k = \sigma$. Its base is $\sigma$ and its slope is at least $\sigma$, so $\mu$ falls in case A. It is not exceptional, since its base $\sigma \leq k$ is smaller than its number of
parts $k + 1$.

\emph{Step 5: $\phi$ is an involution.} If $\lambda$ is in case A with base $s$, then $\mu = \phi(\lambda)$ has slope $s$, and case B for $\mu$ subtracts $1$ from its $s$ largest parts and adds the
part $s$: this is $\lambda$. If $\lambda$ is in case B with slope $\sigma$, then $\mu$ has base $\sigma$, and case A for $\mu$ removes it and adds $1$ to the $\sigma$ largest parts: this is $\lambda$.
So $\phi$ is an involution on the non-exceptional partitions into distinct parts, preserving the size $n$ and changing the number of parts by $\pm 1$. It pairs partitions with an even number of parts
with partitions with an odd number, so the non-exceptional partitions contribute $0$ to $e(n) - o(n)$.

\emph{Step 6: the exceptional partitions.} If $\sigma = k$ and $s = k$, then $\lambda = (2k - 1, 2k - 2, \ldots, k)$ and $n = \sum_{i=k}^{2k-1}i = \frac{k(3k - 1)}{2} = \omega(k)$. If $\sigma = k$
and $s = k + 1$, then $\lambda = (2k, \ldots, k + 1)$ and $n = \frac{k(3k + 1)}{2} = \omega(-k)$. With the empty partition ($k = 0$, $n = 0 = \omega(0)$), every pentagonal number $\omega(j)$ has
exactly one exceptional partition, with $|j|$ parts, and the other numbers have none. So $e(n) - o(n) = (-1)^j$ for $n = \omega(j)$ and $0$ otherwise.
\end{proof}

\begin{example}\label{generating-functions:example-franklin}
For $n = 7$ the partitions into distinct parts are $7$, $6 + 1$, $5 + 2$, $4 + 3$, $4 + 2 + 1$. Franklin's involution pairs $6 + 1 \leftrightarrow 7$ (case A for $6 + 1$: base $1$, slope $1$) and
$4 + 2 + 1 \leftrightarrow 5 + 2$ (case A for $4 + 2 + 1$: base $1$, slope $1$, giving $5 + 2$), and leaves $4 + 3$ exceptional: slope $2$, base $3 = k + 1$. So $e(7) - o(7) = 1$, and indeed
$7 = \omega(-2)$ with sign $(-1)^{-2} = +1$.
\end{example}

\begin{corollary}[Euler's recurrence for $p(n)$]\label{generating-functions:corollary-partition-recurrence}
For $n \geq 1$,
\[
\begin{aligned}
p(n) &= \sum_{j \neq 0}(-1)^{j-1}p\bigl(n - \omega(j)\bigr)\\
&= p(n - 1) + p(n - 2) - p(n - 5) - p(n - 7) + p(n - 12) + \cdots,
\end{aligned}
\]
with $p(m) = 0$ for $m < 0$.
\end{corollary}

\begin{proof}
By Theorem~\ref{generating-functions:theorem-partition-gf} and Lemma~\ref{generating-functions:lemma-product-rules}(2), $\sum_np(n)x^n \cdot \prod_k(1 - x^k) = 1$. Insert the pentagonal number
theorem and compare the coefficients of $x^n$, $n \geq 1$: $\sum_j(-1)^jp(n - \omega(j)) = 0$, and the term $j = 0$ is $p(n)$.
\end{proof}

The recurrence computes $p(n)$ with about $\sqrt n$ additions per value, instead of the $n$ multiplications by $1/(1 - x^k)$. For example $p(10) = p(9) + p(8) - p(5) - p(3) = 30 + 22 - 7 - 3 = 42$.
MacMahon used it to compute $p(n)$ for $n \leq 200$, finding $p(200) = 3\,972\,999\,029\,388$, which Hardy and Ramanujan used to test their asymptotic formula
(Remark~\ref{combinatorics:remark-hardy-ramanujan}).

\begin{corollary}[Euler's recurrence for $\sigma(n)$]\label{generating-functions:corollary-sigma-recurrence}
Let $\sigma(n)$ be the sum of the positive divisors of $n$. For $n \geq 1$,
\[
\begin{aligned}
\sigma(n) &= \sum_{j \neq 0}(-1)^{j-1}\sigma\bigl(n - \omega(j)\bigr)\\
&= \sigma(n - 1) + \sigma(n - 2) - \sigma(n - 5) - \sigma(n - 7) + \cdots,
\end{aligned}
\]
where the terms with $n - \omega(j) < 0$ are omitted, and a term $\sigma(0)$, which occurs exactly when $n$ is pentagonal, is to be read as $n$.
\end{corollary}

\begin{proof}
Let $E = \prod_k(1 - x^k) = \sum_me_mx^m$, so $e_m = (-1)^j$ if $m = \omega(j)$ and $0$ otherwise. By Lemma~\ref{generating-functions:lemma-product-rules}(4),
\[
-x\frac{DE}{E} = \sum_{k \geq 1}\frac{kx^k}{1 - x^k} = \sum_{k \geq 1}\sum_{m \geq 1}kx^{km} = \sum_{n \geq 1}\sigma(n)x^n,
\]
grouping the pairs $(k, m)$ with $km = n$ (Proposition~\ref{generating-functions:proposition-summable}(2)). Write $S = \sum_{n \geq 1}\sigma(n)x^n$; then $-xDE = ES$, and comparing coefficients of
$x^n$: $-ne_n = \sum_{m=0}^{n-1}e_m\sigma(n - m)$. The term $m = 0$ is $\sigma(n)$, so $\sigma(n) = -\sum_{m \geq 1}e_m\sigma(n - m) - ne_n$, which is the statement: the last term is $(-1)^{j-1}n$ when
$n = \omega(j)$.
\end{proof}

For instance $\sigma(12) = \sigma(11) + \sigma(10) - \sigma(7) - \sigma(5) + 12 = 12 + 18 - 8 - 6 + 12 = 28$, which is $1 + 2 + 3 + 4 + 6 + 12$. Euler found this recurrence before he could prove the
pentagonal number theorem, and called it ``a most extraordinary law of the numbers concerning the sums of their divisors''.

\subsection{Jacobi's triple product identity}\label{generating-functions:subsection-triple-product}

The pentagonal number theorem is a special case of an identity in two variables. We work in the ring $\Lambda[[q]]$ of formal power series in $q$ whose coefficients are Laurent polynomials in $z$,
$\Lambda = \ZZ[z, z^{-1}]$.

\begin{theorem}[Jacobi's triple product identity]\label{generating-functions:theorem-triple-product}
In $\Lambda[[q]]$,
\[
\prod_{n \geq 1}(1 - q^{2n})(1 + q^{2n-1}z)(1 + q^{2n-1}z^{-1}) = \sum_{m \in \ZZ}q^{m^2}z^m.
\]
\end{theorem}

\begin{proof}
\emph{Step 1: a finite identity.} We show that for $N \geq 1$, in the ring $\ZZ[q, q^{-1}, z, z^{-1}]$ of Laurent polynomials,
\begin{equation}\label{generating-functions:equation-finite-triple}
\prod_{i=1}^{N}(1 + q^{2i-1}z)(1 + q^{2i-1}z^{-1}) = \sum_{m=-N}^{N}q^{m^2}z^m\binom{2N}{N + m}_{q^2}.
\end{equation}
Each factor $1 + q^{2i-1}z^{-1} = q^{2i-1}z^{-1}(1 + q^{-(2i-1)}z)$, and $\sum_{i=1}^{N}(2i - 1) = N^2$, so the left-hand side is $q^{N^2}z^{-N}\prod_{e}(1 + q^ez)$, the product over the odd
$e$ with $|e| \leq 2N - 1$. These are $e = 2i - (2N - 1)$, $0 \leq i \leq 2N - 1$, so the product is $\prod_{i=0}^{2N-1}(1 + Q^iw)$ with $Q = q^2$ and $w = q^{-(2N-1)}z$. By
Theorem~\ref{generating-functions:theorem-q-binomial} (an identity in $\ZZ[Q, w]$, which we may evaluate in the commutative ring of Laurent polynomials) it equals
$\sum_{j=0}^{2N}Q^{j(j-1)/2}\binom{2N}{j}_Qw^j = \sum_jq^{j(j-1) - (2N-1)j}z^j\binom{2N}{j}_{q^2}$. Multiplying by $q^{N^2}z^{-N}$ and putting $m = j - N$, the exponent of $q$ becomes
$j^2 - 2Nj + N^2 = m^2$, which is \eqref{generating-functions:equation-finite-triple}.

\emph{Step 2: the limit.} Both sides of \eqref{generating-functions:equation-finite-triple} lie in $\Lambda[q] \subseteq \Lambda[[q]]$. As $N \to \infty$, the left-hand side converges to the
infinite product $\prod_{n \geq 1}(1 + q^{2n-1}z)(1 + q^{2n-1}z^{-1})$ (Proposition~\ref{generating-functions:proposition-infinite-product}(3)). On the right, fix an exponent $M$ of $q$: only the
$m$ with $m^2 \leq M$ contribute, and for each such $m$ the coefficient of $q^{M - m^2}$ in $\binom{2N}{N + m}_{q^2}$ counts the partitions of $(M - m^2)/2$ fitting in an $(N + m) \times (N - m)$
box, which for $N - |m| \geq M$ is $p((M - m^2)/2)$ (and $0$ if $M - m^2$ is odd). So the right-hand side converges to $\sum_mq^{m^2}z^m\sum_kp(k)q^{2k} = \sum_mq^{m^2}z^m\prod_{n \geq 1}(1 - q^{2n})^{-1}$,
by Theorem~\ref{generating-functions:theorem-partition-gf} with $x = q^2$ (Lemma~\ref{generating-functions:lemma-product-rules}(5)). Multiplying by $\prod_n(1 - q^{2n})$
(Lemma~\ref{generating-functions:lemma-product-rules}(1), (2)) gives the theorem.
\end{proof}

\begin{corollary}\label{generating-functions:corollary-triple-product}
In $\ZZ[[q]]$:
\begin{enumerate}
\item (Gauss) $\displaystyle\sum_{m \in \ZZ}(-1)^mq^{m^2} = \prod_{n \geq 1}\frac{1 - q^n}{1 + q^n} = 1 - 2q + 2q^4 - 2q^9 + \cdots$;
\item $\displaystyle\sum_{m \in \ZZ}q^{m^2} = \prod_{n \geq 1}(1 - q^{2n})(1 + q^{2n-1})^2$;
\item the pentagonal number theorem again: $\prod_{n \geq 1}(1 - x^n) = \sum_j(-1)^jx^{j(3j-1)/2}$.
\end{enumerate}
\end{corollary}

\begin{proof}
(1) Apply the coefficient homomorphism $\Lambda \to \ZZ$, $z \mapsto -1$ (Lemma~\ref{generating-functions:lemma-product-rules}(5)): $\sum_m(-1)^mq^{m^2} = \prod_n(1 - q^{2n})(1 - q^{2n-1})^2$. Now
$\prod_n(1 - q^{2n})(1 - q^{2n-1}) = \prod_n(1 - q^n)$, splitting by parity (Lemma~\ref{generating-functions:lemma-product-rules}(3)), and
$\prod_n(1 - q^{2n-1}) = \prod_n\frac{1}{1 + q^n}$ by Corollary~\ref{generating-functions:corollary-glaisher} with $d = 2$ (in the form $\prod_n(1 + q^n)\prod_{n \text{ odd}}(1 - q^n) = 1$, which
follows from its proof with $x = q$, after inverting both sides). Multiplying gives (1).
(2) Apply $z \mapsto 1$.
(3) Replace $q$ by $q^3$ and then $z$ by $-q$, keeping track of the exponents. Since the substitution $z \mapsto -q$ does not preserve $\Lambda[[q]]$, we argue on the finite identity
\eqref{generating-functions:equation-finite-triple} instead: with $q^3$ for $q$ and $-q$ for $z$, an identity of Laurent polynomials, the left-hand side becomes
$\prod_{i=1}^{N}(1 - q^{6i-2})(1 - q^{6i-4})$ and the right-hand side $\sum_{m=-N}^{N}(-1)^mq^{3m^2+m}\binom{2N}{N + m}_{q^6}$, both polynomials in $q$. Letting $N \to \infty$ as in Step 2 of the proof,
and multiplying by $\prod_n(1 - q^{6n})$, gives $\prod_n(1 - q^{6n})(1 - q^{6n-2})(1 - q^{6n-4}) = \sum_m(-1)^mq^{3m^2+m}$. The left-hand side is $\prod_n(1 - q^{2n})$, splitting the even exponents
modulo $6$. Now replace $q^2$ by $x$: the exponent $3m^2 + m$ is even, and with $j = -m$, $(3m^2 + m)/2 = j(3j - 1)/2$. Formally: both sides lie in the image of the injective substitution
$f \mapsto f(q^2)$, so the identity holds for the preimages.
\end{proof}

\begin{remark}\label{generating-functions:remark-partition-identities}
The triple product identity is the source of a great many identities. A limiting form of it gives Jacobi's identity
$\prod_{n \geq 1}(1 - x^n)^3 = \sum_{m \geq 0}(-1)^m(2m + 1)x^{m(m+1)/2}$; this and the \emph{Rogers--Ramanujan identities}
\[
\begin{aligned}
\sum_{k \geq 0}\frac{x^{k^2}}{(x)_k} &= \prod_{n \geq 0}\frac{1}{(1 - x^{5n+1})(1 - x^{5n+4})},\\
\sum_{k \geq 0}\frac{x^{k^2+k}}{(x)_k} &= \prod_{n \geq 0}\frac{1}{(1 - x^{5n+2})(1 - x^{5n+3})},
\end{aligned}
\]
are proved in \cite{andrews-partitions}. The first Rogers--Ramanujan identity says that the partitions of $n$ whose parts differ pairwise by at least $2$ are equinumerous with the partitions of $n$
into parts congruent to $\pm 1$ modulo $5$; bijective proofs exist, but they are complicated. Ramanujan also discovered the congruences $p(5n + 4) \equiv 0 \pmod 5$, $p(7n + 5) \equiv 0 \pmod 7$ and
$p(11n + 6) \equiv 0 \pmod{11}$, whose proofs use the theory of modular forms (Chapter~\ref{modular-forms}).
\end{remark}

\section{Exponential generating functions}\label{generating-functions:section-exponential}

When the objects being counted are built on a set of \emph{labelled} elements, such as permutations, set partitions or graphs on $[n]$, the natural generating function divides the $n$-th term by
$n!$. The reason is that the product of two such series then counts pairs of structures placed on complementary subsets, with the binomial coefficient choosing the subsets. In this section $R$
contains $\QQ$.

\subsection{Labelled structures}\label{generating-functions:subsection-labelled}

\begin{definition}\label{generating-functions:definition-egf}
The \emph{exponential generating function} of a sequence $(a_n)_{n \geq 0}$ is $\sum_na_n\frac{x^n}{n!} \in R[[x]]$. For a series $f$ we write $n![x^n]f$ for the coefficient of $x^n/n!$ in $f$.
\end{definition}

\begin{definition}\label{generating-functions:definition-species}
A \emph{class of labelled structures} $\mathcal{F}$ assigns to every finite set $S$ a finite set $\mathcal{F}[S]$, the structures \emph{on} $S$, and to every bijection $\tau \colon S \to T$ a bijection
$\mathcal{F}[\tau] \colon \mathcal{F}[S] \to \mathcal{F}[T]$, the \emph{relabelling} along $\tau$, such that $\mathcal{F}[\id_S] = \id$ and $\mathcal{F}[\tau \circ \rho] = \mathcal{F}[\tau] \circ \mathcal{F}[\rho]$.
Its \emph{counting sequence} is $f_n = |\mathcal{F}[[n]]|$, which is also $|\mathcal{F}[S]|$ for every $S$ with $n$ elements, and its exponential generating function is $F(x) = \sum_nf_nx^n/n!$.
\end{definition}

Examples: the permutations of $S$ (relabelled by conjugation, $\sigma \mapsto \tau\sigma\tau^{-1}$); the linear orders of $S$; the set partitions of $S$; the graphs with vertex set $S$, that is, the
subsets of $\binom{S}{2}$; the class $\mathcal{E}$ with exactly one structure on every set (the set $S$ itself), with exponential generating function $\sum_nx^n/n! = e^x$; and the class $\mathcal{X}$ of
singletons, with one structure on sets with one element and none otherwise, with generating function $x$. Permutations and linear orders both have $n!$ structures on an $n$-set, and generating
function $\frac{1}{1 - x}$; they are different classes, since relabelling acts differently, but only the counting sequence matters here. The functoriality is what makes $|\mathcal{F}[S]|$ depend only on
$|S|$, and it will let us transport structures from a block of a partition to a standard set.

\begin{theorem}[Labelled product]\label{generating-functions:theorem-labelled-product}
For classes $\mathcal{F}$ and $\mathcal{G}$, let $(\mathcal{F} \cdot \mathcal{G})[S] = \bigsqcup_{(A, B)}\mathcal{F}[A] \times \mathcal{G}[B]$, the disjoint union over the ordered pairs of disjoint
sets with $A \cup B = S$. Its counting sequence is $h_n = \sum_{k=0}^{n}\binom{n}{k}f_kg_{n-k}$, and its exponential generating function is $F(x)G(x)$. More generally, the class of $k$-tuples of
structures on the blocks of ordered decompositions $S = A_1 \sqcup \cdots \sqcup A_k$ has the generating function $F_1(x) \cdots F_k(x)$, and the number of such tuples on $[n]$ with $|A_i| = n_i$ is
$\binom{n}{n_1, \ldots, n_k}f^{(1)}_{n_1} \cdots f^{(k)}_{n_k}$.
\end{theorem}

\begin{proof}
For $S = [n]$, choose $A$ with $k$ elements ($\binom{n}{k}$ ways), a structure on $A$ ($f_k$ ways) and one on $B = [n] \setminus A$ ($g_{n-k}$ ways); the rules of sum and product give $h_n$. On the
other side, $n![x^n](FG) = n!\sum_{k}\frac{f_k}{k!}\frac{g_{n-k}}{(n - k)!} = \sum_k\binom{n}{k}f_kg_{n-k}$. The case of $k$ factors follows by induction, or directly with the multinomial
coefficient counting the ordered decompositions with prescribed sizes (Definition~\ref{combinatorics:definition-multinomial}).
\end{proof}

\begin{example}\label{generating-functions:example-labelled-product}
\begin{enumerate}
\item The pairs (subset $A$, the rest) are the structures of $\mathcal{E} \cdot \mathcal{E}$, with generating function $e^xe^x = e^{2x}$: there are $2^n$ of them.
\item A permutation of $S$ together with a choice of its fixed points is the same as a set $A$ of fixed points and a derangement of $S \setminus A$. So $\frac{1}{1 - x} = e^xD(x)$, with $D$ the
exponential generating function of the derangements, and $D(x) = \frac{e^{-x}}{1 - x}$: comparing coefficients, $D_n = n!\sum_{j \leq n}\frac{(-1)^j}{j!}$, which is
Theorem~\ref{combinatorics:theorem-derangements} again.
\end{enumerate}
\end{example}

\subsection{The exponential formula}\label{generating-functions:subsection-exponential-formula}

\begin{theorem}[Exponential formula]\label{generating-functions:theorem-exponential-formula}
Let $\mathcal{C}$ be a class of labelled structures with $\mathcal{C}[\emptyset] = \emptyset$ (``connected'' structures) and generating function $C(x)$. Let $\mathcal{A}[S]$ be the set of pairs
$(\pi, (c_B)_{B \in \pi})$ of a set partition $\pi$ of $S$ and a structure $c_B \in \mathcal{C}[B]$ on each block (``assemblies'' of connected structures). Then the number of assemblies on $[n]$
with exactly $k$ blocks has exponential generating function $C(x)^k/k!$, and in $R[y][[x]]$,
\[
\begin{aligned}
&\sum_{n,k}\bigl|\{\text{assemblies on } [n] \text{ with } k \text{ blocks}\}\bigr|\,y^k\frac{x^n}{n!} = \exp(yC(x)),\\
&\sum_n|\mathcal{A}[[n]]|\frac{x^n}{n!} = \exp(C(x)).
\end{aligned}
\]
\end{theorem}

\begin{proof}
Count the pairs (assembly with $k$ blocks, numbering of its blocks). A numbering turns the set partition into an ordered decomposition $[n] = A_1 \sqcup \cdots \sqcup A_k$ into nonempty blocks, and
the assembly into a $k$-tuple of structures $c_i \in \mathcal{C}[A_i]$; conversely such a tuple determines the assembly and the numbering. Since the blocks are nonempty and distinct, each assembly
has exactly $k!$ numberings. So, with $a_{n,k}$ the number of assemblies,
\[
k!\,a_{n,k} = \bigl|\{(A_1, \ldots, A_k, c_1, \ldots, c_k)\}\bigr| = n![x^n]C(x)^k,
\]
by Theorem~\ref{generating-functions:theorem-labelled-product} (the condition $A_i \neq \emptyset$ is automatic, as $\mathcal{C}[\emptyset] = \emptyset$). Hence
$\sum_na_{n,k}x^n/n! = C^k/k!$, and summing over $k$ (a summable family, as $\operatorname{ord}(C^k) \geq k$) gives $\sum_ky^kC^k/k! = \exp(yC)$, where $\exp(yC)$ is the composition of $\exp$ with
$yC$, which has constant term $0$. Setting $y = 1$ (Lemma~\ref{generating-functions:lemma-product-rules}(5)) gives the second formula.
\end{proof}

The theorem says: \emph{to pass from connected structures to all structures, exponentiate}; and conversely, $C = \log(A)$ recovers the connected structures from all structures. The condition
$\mathcal{C}[\emptyset] = \emptyset$ is necessary for $\exp(C)$ to be defined, and natural: the empty set has one assembly, the empty partition, which accounts for the constant term $1$.

\begin{example}[Set partitions]\label{generating-functions:example-egf-partitions}
Take $\mathcal{C}$ with one structure on every nonempty set: $C = e^x - 1$. Assemblies are set partitions, counted by the Stirling numbers $S(n, k)$ and the Bell numbers $B_n$
(Section~\ref{combinatorics:section-stirling}):
\[
\begin{aligned}
&\sum_{n,k}S(n, k)y^k\frac{x^n}{n!} = \exp\bigl(y(e^x - 1)\bigr),\\
&\sum_nS(n, k)\frac{x^n}{n!} = \frac{(e^x - 1)^k}{k!}, \qquad \sum_nB_n\frac{x^n}{n!} = e^{e^x - 1}.
\end{aligned}
\]
The derivative of the last series is $e^x\exp(e^x - 1)$ (Proposition~\ref{generating-functions:proposition-exp-log}(5)), and comparing coefficients of $x^n/n!$ by
Theorem~\ref{generating-functions:theorem-labelled-product} gives $B_{n+1} = \sum_k\binom{n}{k}B_k$, the recurrence of Proposition~\ref{combinatorics:proposition-bell}. Expanding
$(e^x - 1)^k = \sum_j\binom{k}{j}(-1)^{k-j}e^{jx}$ gives back the formula $S(n, k) = \frac{1}{k!}\sum_j(-1)^{k-j}\binom{k}{j}j^n$.
\end{example}

\begin{example}[Permutations and their cycles]\label{generating-functions:example-egf-cycles}
A permutation is an assembly of cycles, counting each fixed point as a cycle of length $1$ (Proposition~\ref{groups:proposition-cycle-decomposition}): the blocks are the orbits, and on each
orbit the permutation acts as a cyclic permutation. On a set with $k$ elements there are $(k - 1)!$ cyclic permutations, so
$C(x) = \sum_{k \geq 1}\frac{(k-1)!}{k!}x^k = \sum_k\frac{x^k}{k} = -\log(1 - x)$ (Definition~\ref{generating-functions:definition-exp-log}, with $x$ replaced by $-x$). Then
$\exp(C) = \exp(\log(1 - x))^{-1} = \frac{1}{1 - x}$, the $n!$ permutations, by Proposition~\ref{generating-functions:proposition-exp-log}(2) and (3); and with $y$ counting cycles,
\[
\exp(-y\log(1 - x)) = (1 - x)^{-y} = \sum_n\binom{-y}{n}(-x)^n = \sum_ny^{\overline{n}}\frac{x^n}{n!}
\]
by Theorem~\ref{generating-functions:theorem-binomial-series} in $\QQ[y][[x]]$ and upper negation (Proposition~\ref{combinatorics:proposition-binomial-polynomial}(3)). So the number of permutations of $[n]$ with $k$
cycles is the coefficient of $y^k$ in $y(y + 1) \cdots (y + n - 1)$: this is Theorem~\ref{combinatorics:theorem-stirling-first}(2). Restricting the allowed cycle lengths to a set $L$ gives
$\exp(\sum_{k \in L}x^k/k)$:
\begin{enumerate}
\item derangements ($L = \{2, 3, \ldots\}$): $\exp(-\log(1 - x) - x) = \frac{e^{-x}}{1 - x}$;
\item involutions ($L = \{1, 2\}$): $\exp(x + x^2/2)$; differentiating, $I^{\prime} = (1 + x)I$, so $I_{n+1} = I_n + nI_{n-1}$ (Exercise~\ref{combinatorics:exercise-involutions});
\item permutations with $\sigma^m = \id$ ($L$ the divisors of $m$): $\exp(\sum_{d \mid m}x^d/d)$.
\end{enumerate}
\end{example}

\begin{example}[Connected graphs]\label{generating-functions:example-connected-graphs}
A graph on $S$ is a set $E$ of $2$-element subsets of $S$ (its edges). The vertices of $S$ split into the classes of the equivalence relation generated by ``$\{u, v\} \in E$'', its \emph{components};
the graph is \emph{connected} if $S \neq \emptyset$ and there is one class. Every edge lies inside a component, and the graph induced on a component is connected, since two vertices of a class are
joined by a chain of edges whose vertices all lie in the class. So a graph is the same as a set partition of $S$ (into components) with a connected graph on each block: graphs are assemblies of
connected graphs. There are $2^{\binom{n}{2}}$ graphs on $[n]$, so the exponential generating function $C$ of connected graphs is
\[
C(x) = \log\Bigl(\sum_{n \geq 0}2^{\binom{n}{2}}\frac{x^n}{n!}\Bigr),
\]
a formal identity although the series inside the logarithm converges nowhere except at $0$. Writing $G$ for that series, $DG = G \cdot DC$, which gives the recurrence
$c_n = g_n - \frac{1}{n}\sum_{k=1}^{n-1}k\binom{n}{k}c_kg_{n-k}$ with $g_m = 2^{\binom{m}{2}}$, and the values $c_1, \ldots, c_6 = 1, 1, 4, 38, 728, 26704$: already for $n = 6$ about $81\%$ of the $32768$ graphs are connected.
\end{example}

\subsection{Bernoulli numbers and sums of powers}\label{generating-functions:subsection-bernoulli}

\begin{definition}\label{generating-functions:definition-bernoulli}
The \emph{Bernoulli numbers} $B_n \in \QQ$ are defined by $\frac{x}{e^x - 1} = \sum_{n \geq 0}B_n\frac{x^n}{n!}$, where $\frac{e^x - 1}{x} = \sum_n\frac{x^n}{(n + 1)!}$ has constant term $1$, so its
inverse exists in $\QQ[[x]]$. (They are unrelated to the Bell numbers, which unfortunately share the letter.)
\end{definition}

\begin{proposition}\label{generating-functions:proposition-bernoulli}
\begin{enumerate}
\item $B_0 = 1$, and $\sum_{j=0}^{n}\binom{n + 1}{j}B_j = 0$ for $n \geq 1$; so $B_1 = -\frac12$, $B_2 = \frac16$, $B_3 = 0$, $B_4 = -\frac1{30}$, $B_5 = 0$, $B_6 = \frac1{42}$.
\item $B_n = 0$ for every odd $n \geq 3$.
\item (Faulhaber) For $m \geq 1$ and $n \geq 0$, $\displaystyle\sum_{k=0}^{m-1}k^n = \frac{1}{n + 1}\sum_{j=0}^{n}\binom{n + 1}{j}B_jm^{n+1-j}$.
\end{enumerate}
\end{proposition}

\begin{proof}
(1) Multiply $\sum_jB_jx^j/j!$ by $\frac{e^x - 1}{x} = \sum_i\frac{x^i}{(i+1)!}$; the product is $1$, so for $n \geq 1$, $\sum_{j=0}^{n}\frac{B_j}{j!\,(n - j + 1)!} = 0$; multiply by $(n + 1)!$.
(2) Let $f = \frac{x}{e^x - 1} + \frac{x}{2}$. Substituting $-x$ for $x$ (Lemma~\ref{generating-functions:lemma-composition}), and using $e^{-x} = 1/e^x$,
$\frac{-x}{e^{-x} - 1} = \frac{xe^x}{e^x - 1} = \frac{x}{e^x - 1} + x$, so $f(-x) = \frac{x}{e^x - 1} + x - \frac{x}{2} = f(x)$. An even series has no odd coefficients, and $f$ has the
coefficients $B_n/n!$ for $n \geq 2$.
(3) By Proposition~\ref{generating-functions:proposition-exp-log}(2), $(e^x - 1)\sum_{k=0}^{m-1}e^{kx} = e^{mx} - 1$, a telescoping sum. So
\[
\begin{aligned}
\sum_{n}\Bigl(\sum_{k=0}^{m-1}k^n\Bigr)\frac{x^n}{n!} &= \sum_{k=0}^{m-1}e^{kx} = \frac{x}{e^x - 1} \cdot \frac{e^{mx} - 1}{x}\\
&= \Bigl(\sum_jB_j\frac{x^j}{j!}\Bigr)\Bigl(\sum_{i \geq 0}\frac{m^{i+1}x^i}{(i + 1)!}\Bigr).
\end{aligned}
\]
The coefficient of $x^n/n!$ on the right is $\sum_{j}\binom{n}{j}B_j\frac{m^{n-j+1}}{n - j + 1}$, and $\binom{n}{j}\frac{1}{n - j + 1} = \frac{1}{n + 1}\binom{n + 1}{j}$.
\end{proof}

For $n = 2$: $\sum_{k<m}k^2 = \frac13(m^3 - \frac32m^2 + \frac12m) = \frac{m(m - 1)(2m - 1)}{6}$. The Bernoulli numbers reappear in the values $\zeta(2k)$ of the Riemann zeta function, in the
Euler--Maclaurin formula, and in Kummer's work on Fermat's last theorem.

\section{Lagrange inversion}\label{generating-functions:section-lagrange}

Many generating functions are defined implicitly, by an equation $T = x\phi(T)$ expressing that a structure is a root with a collection of substructures of the same kind attached: binary trees,
plane trees, rooted labelled trees. Lagrange's formula gives the coefficients of $T$ directly. In this section $R$ contains $\QQ$.

\subsection{Laurent series and residues}\label{generating-functions:subsection-laurent}

\begin{definition}\label{generating-functions:definition-laurent}
A \emph{formal Laurent series} over $R$ is a family $f = (f_n)_{n \in \ZZ}$ of elements of $R$ with $f_n = 0$ for all sufficiently negative $n$, written $\sum_nf_nx^n$. They form a ring $R((x))$ under
termwise addition and the product $(fg)_n = \sum_{i+j=n}f_ig_j$, a finite sum. The \emph{residue} of $f$ is $\operatorname{res}(f) = f_{-1}$, and the derivative is $Df = \sum_nnf_nx^{n-1}$.
\end{definition}

The ring axioms, the Leibniz rule and the rules for summable families are proved as for $R[[x]]$; every Laurent series is $x^{-m}g$ with $m \geq 0$ and $g \in R[[x]]$, and $R[[x]]$ is a subring. If
$K$ is a field, $K((x))$ is a field, the field of fractions of $K[[x]]$: a nonzero $f$ is $x^mu$ with $m \in \ZZ$ and $u$ a unit of $K[[x]]$ (Proposition~\ref{generating-functions:proposition-field-coefficients}),
with inverse $x^{-m}u^{-1}$.

\begin{lemma}\label{generating-functions:lemma-residues}
Let $f \in R((x))$.
\begin{enumerate}
\item $\operatorname{res}(Df) = 0$.
\item If $f = x^mu$ with $m \in \ZZ$ and $u \in R[[x]]$ with $u(0)$ a unit, then $f$ is a unit of $R((x))$ and $\operatorname{res}(Df/f) = m$.
\item Let $g \in R[[x]]$ with $g(0) = 0$ and $g_1$ a unit. The composition $f \mapsto f(g) = \sum_nf_ng^n$ is defined on $R((x))$, where $g^{-1} \in R((x))$ by (2); it is a ring homomorphism, and
\[
\operatorname{res}\bigl(f(g) \cdot Dg\bigr) = \operatorname{res}(f).
\]
\end{enumerate}
\end{lemma}

\begin{proof}
(1) The coefficient of $x^{-1}$ in $Df$ is $0 \cdot f_0 = 0$.
(2) $u$ is a unit of $R[[x]]$ (Proposition~\ref{rings:proposition-power-series-units}), so $f^{-1} = x^{-m}u^{-1}$. By the Leibniz rule, $Df = mx^{m-1}u + x^mDu$, so $Df/f = \frac{m}{x} + \frac{Du}{u}$,
and $Du/u \in R[[x]]$ has no residue.
(3) Write $g = xv$ with $v(0) = g_1$ a unit; then $g^{-1} = x^{-1}v^{-1}$ and $g^n$ has order exactly $n$ for all $n \in \ZZ$, with lowest coefficient $g_1^n$. The family $(f_ng^n)$ is summable, since
$f_n = 0$ for $n$ below some $n_0$ and $\operatorname{ord}(g^n) = n$; so $f(g)$ is defined, and it is a homomorphism by the computation of Lemma~\ref{generating-functions:lemma-composition}(1).
For the residue formula, both sides are additive in $f$ and only the terms with $n_0 \leq n \leq -1$ matter: for $n \geq 0$, $f_ng^nDg \in R[[x]]$ and $f_nx^n$ have residue $0$. So it suffices
to treat $f = x^n$ with $n < 0$. If $n \neq -1$, then $g^nDg = D(g^{n+1})/(n + 1)$ has residue $0$ by (1), like $x^n$. If $n = -1$, then $g^{-1}Dg = Dg/g$ has residue $1$ by (2), like $x^{-1}$.
\end{proof}

\subsection{The inversion formula}\label{generating-functions:subsection-lagrange-formula}

\begin{theorem}[Lagrange inversion]\label{generating-functions:theorem-lagrange}
Let $\phi \in R[[t]]$ with $\phi(0)$ a unit. There is a unique $T \in R[[x]]$ with $T(0) = 0$ and $T = x\phi(T)$. For every $H \in R[[t]]$ and $n \geq 1$,
\[
[x^n]H(T) = \frac{1}{n}[t^{n-1}]\,H^{\prime}(t)\phi(t)^n,
\]
and in particular
\[
[x^n]T^k = \frac{k}{n}[t^{n-k}]\,\phi(t)^n \quad (k \geq 1).
\]
\end{theorem}

\begin{proof}
\emph{Existence and uniqueness.} Since $\phi$ is a unit, $T = x\phi(T)$ is equivalent to $F(T) = x$ with $F(t) = t/\phi(t)$. Here $F(0) = 0$ and $F_1 = \phi(0)^{-1}$ is a unit, so by
Proposition~\ref{generating-functions:proposition-compositional-inverse} there is exactly one $T$ with $T(0) = 0$ and $F \circ T = x$, namely the compositional inverse of $F$; it has $T_1 = \phi(0)$,
a unit.
\emph{The formula.} For $n \geq 1$, $n[x^n]H(T) = [x^{n-1}]D(H(T)) = \operatorname{res}\bigl(x^{-n}D(H(T))\bigr)$. By the chain rule (Lemma~\ref{generating-functions:lemma-derivative}(3)), $D(H(T)) = H^{\prime}(T)DT$. Since $x = F(T)$ and composition with $T$
is a homomorphism on Laurent series (Lemma~\ref{generating-functions:lemma-residues}(3)), $x^{-n} = (F^{-n})(T)$, where $F(t)^{-n} = t^{-n}\phi(t)^n \in R((t))$. Hence
\[
\begin{aligned}
n[x^n]H(T) &= \operatorname{res}\Bigl(\bigl(F^{-n}H^{\prime}\bigr)(T) \cdot DT\Bigr)\\
&= \operatorname{res}\bigl(t^{-n}\phi(t)^nH^{\prime}(t)\bigr) = [t^{n-1}]\,\phi(t)^nH^{\prime}(t),
\end{aligned}
\]
by Lemma~\ref{generating-functions:lemma-residues}(3) with $g = T$. For $H = t^k$, $H^{\prime} = kt^{k-1}$ and $[t^{n-1}]kt^{k-1}\phi^n = k[t^{n-k}]\phi^n$.
\end{proof}

\begin{example}[Catalan numbers again]\label{generating-functions:example-lagrange-catalan}
The series $T = xB$, where $B = 1 + xB^2$ is the generating function of binary trees (Theorem~\ref{generating-functions:theorem-catalan}), satisfies $T = x + T^2$, that is, $T = \frac{x}{1 - T} = x\phi(T)$ with $\phi(t) = (1 - t)^{-1}$.
By Lagrange and Example~\ref{generating-functions:example-geometric}(3), $C_{n-1} = [x^n]T = \frac{1}{n}[t^{n-1}](1 - t)^{-n} = \frac{1}{n}\binom{2n - 2}{n - 1}$. More generally
$[x^{n+k}]T^k = \frac{k}{n + k}\binom{2n + k - 1}{n}$ counts the $k$-tuples of binary trees with $n$ nodes in all (Exercise~\ref{generating-functions:exercise-ballot-numbers}).
\end{example}

\begin{definition}\label{generating-functions:definition-mary-tree}
For $m \geq 1$, the sets $\mathcal{T}^{(m)}_n$ of \emph{$m$-ary trees with $n$ nodes} are defined by $\mathcal{T}^{(m)}_0 = \{\varepsilon\}$ and
$\mathcal{T}^{(m)}_n = \bigsqcup_{n_1 + \cdots + n_m = n - 1}\mathcal{T}^{(m)}_{n_1} \times \cdots \times \mathcal{T}^{(m)}_{n_m}$: a nonempty $m$-ary tree is a root with an ordered $m$-tuple of
$m$-ary trees. For $m = 2$ these are the binary trees of Definition~\ref{generating-functions:definition-binary-tree}.
\end{definition}

\begin{corollary}[Fuss--Catalan numbers]\label{generating-functions:corollary-fuss-catalan}
The number of $m$-ary trees with $n$ nodes is $\frac{1}{(m - 1)n + 1}\binom{mn}{n}$.
\end{corollary}

\begin{proof}
The generating function $A$ satisfies $A = 1 + xA^m$, by the rule of product, so $T = A - 1$ satisfies $T = x(1 + T)^m$, and $T(0) = 0$. By Lagrange with $\phi = (1 + t)^m$, for $n \geq 1$,
$[x^n]T = \frac{1}{n}[t^{n-1}](1 + t)^{mn} = \frac{1}{n}\binom{mn}{n - 1}$, and $\frac{1}{n}\binom{mn}{n - 1} = \frac{(mn)!}{n!\,(mn - n + 1)!} = \frac{1}{(m - 1)n + 1}\binom{mn}{n}$. For $n = 0$ both
sides are $1$.
\end{proof}

For $m = 3$ the numbers are $1, 1, 3, 12, 55, 273, \ldots$, and for $m = 1$ they are all $1$ (a $1$-ary tree is a path).

\subsection{The cycle lemma}\label{generating-functions:subsection-cycle-lemma}

The factor $\frac{1}{(m - 1)n + 1}$ in the Fuss--Catalan numbers, like the factor $\frac{1}{n + 1}$ in the Catalan numbers, has a combinatorial explanation: among the cyclic shifts of a sequence,
exactly one has a positivity property.

\begin{lemma}[Cycle lemma]\label{generating-functions:lemma-cycle}
Let $a_1, \ldots, a_N$ be integers with $a_1 + \cdots + a_N = 1$. Exactly one of the $N$ cyclic shifts $(a_{j+1}, \ldots, a_N, a_1, \ldots, a_j)$, $0 \leq j < N$, has all its partial sums positive;
and these $N$ shifts are pairwise distinct sequences.
\end{lemma}

\begin{proof}
Put $s_0 = 0$ and $s_i = a_1 + \cdots + a_i$, and extend by $s_{i+N} = s_i + 1$ for all $i \geq 0$, which matches the partial sums of the periodic sequence $a_{i+N} = a_i$. The partial sums of the
shift starting after position $j$ are $s_{j+t} - s_j$, $1 \leq t \leq N$. Let $j$ be the \emph{last} index in $\{0, \ldots, N - 1\}$ at which $s_j = \min(s_0, \ldots, s_{N-1})$.
\emph{Existence.} For $j < i \leq N - 1$, $s_i > s_j$ because $j$ is the last minimum; for $N \leq i < j + N$, $s_i = s_{i-N} + 1 \geq s_j + 1 > s_j$; and $s_{j+N} = s_j + 1$. So all partial sums of the
shift after $j$ are positive.
\emph{Uniqueness.} If the shift after $k$ has positive partial sums, then $s_i > s_k$ for $k < i \leq N - 1$, and $s_i + 1 = s_{i+N} > s_k$, that is $s_i \geq s_k$, for $0 \leq i < k$. So $s_k$ is the
minimum and $k$ is its last occurrence: $k = j$.
\emph{Distinctness.} If two shifts coincide, the sequence is invariant under a shift by some $d$ with $0 < d < N$, hence under the subgroup of $\ZZ/N\ZZ$ generated by $d$, of order $N/e$ with
$e = \gcd(d, N) < N$; then the sequence consists of $N/e$ repetitions of its first $e$ terms, and its sum $1$ is divisible by $N/e > 1$, which is absurd.
\end{proof}

\begin{corollary}\label{generating-functions:corollary-cycle-lemma-count}
Let $m \geq 1$ and $n \geq 0$. The number of sequences consisting of $n$ terms equal to $1 - m$ and $(m - 1)n + 1$ terms equal to $1$, all of whose partial sums are positive, is
$\frac{1}{mn + 1}\binom{mn + 1}{n} = \frac{1}{(m - 1)n + 1}\binom{mn}{n}$.
\end{corollary}

\begin{proof}
There are $N = mn + 1$ terms with sum $(m - 1)n + 1 - (m - 1)n = 1$, and $\binom{N}{n}$ such sequences. The cyclic group $\ZZ/N\ZZ$ acts on them by shifts; by the cycle lemma each orbit has $N$
elements and contains exactly one sequence with positive partial sums. So the number of those is $\binom{N}{n}/N$, and $\frac{1}{mn + 1}\binom{mn + 1}{n} = \frac{(mn)!}{n!\,((m - 1)n + 1)!}$.
\end{proof}

For $m = 2$ the sequences are $+1$ followed by a Dyck word (with $U = +1$, $D = -1$): their number is $C_n$, a third proof of the formula for the Catalan numbers, and the shortest. The two counts of
Corollaries~\ref{generating-functions:corollary-fuss-catalan} and~\ref{generating-functions:corollary-cycle-lemma-count} agree, and indeed reading an $m$-ary tree in preorder (root first, then
the subtrees from left to right), writing $1 - m$ for each node and $+1$ for each empty subtree, and reversing the resulting word, gives a bijection; we do not need it.

\subsection{Labelled trees and Cayley's formula}\label{generating-functions:subsection-cayley}

\begin{definition}\label{generating-functions:definition-rooted-tree}
A \emph{rooted forest} on a finite set $S$ is a pair $(\mathcal{R}, p)$ of a subset $\mathcal{R} \subseteq S$, the \emph{roots}, and a map $p \colon S \setminus \mathcal{R} \to S$, the \emph{parent
map}, such that for every $v \in S$ some iterate $p^k(v)$, $k \geq 0$, lies in $\mathcal{R}$. A \emph{rooted tree} is a rooted forest with one root. The \emph{edges} of $(\mathcal{R}, p)$ are the
sets $\{v, p(v)\}$, $v \in S \setminus \mathcal{R}$; a \emph{tree} on $S$ is a set of edges that arises from some rooted tree on $S$.
\end{definition}

The iterates $v, p(v), p^2(v), \ldots$ are the path from $v$ up to its root; the condition excludes cycles, since a cycle $v = p^k(v)$ would never reach a root. Relabelling along a bijection
transports rooted trees, so these are classes of labelled structures.

\begin{lemma}\label{generating-functions:lemma-trees-roots}
Let $S$ be a finite set with $n \geq 1$ elements.
\begin{enumerate}
\item For a rooted forest $(\mathcal{R}, p)$, the map $v \mapsto \{v, p(v)\}$ is injective on $S \setminus \mathcal{R}$; so a rooted tree has $n - 1$ edges.
\item A tree on $S$ together with a vertex $r \in S$ arises from exactly one rooted tree with root $r$.
\end{enumerate}
Consequently the number of rooted trees on $S$ is $n$ times the number of trees on $S$.
\end{lemma}

\begin{proof}
(1) If $\{v, p(v)\} = \{w, p(w)\}$ with $v \neq w$, then $w = p(v)$ and $v = p(w)$, so $v = p^2(v)$ and the iterates of $v$ alternate between $v$ and $w$, two non-roots: they never reach a root.
\emph{(2), uniqueness.} Let $p$ and $p^{\prime}$ be parent maps with root $r$ and the same edge set $E$. Let $d(v)$ be the least $k$ with $p^k(v) = r$. We show $p^{\prime}(v) = p(v)$ by induction on
$d(v) \geq 1$. By (1) for $p^{\prime}$, the edge $\{v, p(v)\}$ is $\{u, p^{\prime}(u)\}$ for exactly one $u \neq r$, and $u \in \{v, p(v)\}$. If $u = v$ we are done. Otherwise $u = p(v) \neq r$ and
$p^{\prime}(p(v)) = v$; but $d(p(v)) = d(v) - 1$, so by induction $p^{\prime}(p(v)) = p(p(v))$, and $p(p(v)) = v$ is impossible, as in (1).
\emph{(2), existence.} Let $E$ come from a rooted tree $p$ with root $r_0$, and let $r \in S$. Let $r = v_0, v_1 = p(v_0), \ldots, v_d = r_0$ be the path from $r$ to $r_0$. Define $p^{\prime}$ on
$S \setminus \{r\}$ by $p^{\prime}(v_{i+1}) = v_i$ for $0 \leq i < d$, and $p^{\prime}(u) = p(u)$ for $u$ not on the path. It has the same edges, the path being reversed. Every $u$ reaches $r$: its
iterates under $p^{\prime}$ agree with those under $p$ until they first meet the path (which they do, at $r_0$ at the latest), at some $v_i$, and from there $p^{\prime}$ descends $v_i, v_{i-1}, \ldots, v_0 = r$.
The consequence: by (2) the map (rooted tree) $\mapsto$ (its edge set, its root) is a bijection onto the pairs (tree, vertex).
\end{proof}

\begin{theorem}[Cayley's formula]\label{generating-functions:theorem-cayley}
For $n \geq 1$ there are $n^{n-1}$ rooted trees and $n^{n-2}$ trees on $[n]$.
\end{theorem}

\begin{proof}[First proof, by Lagrange inversion]
Removing the root $r$ of a rooted tree on $S$ leaves the rooted forest on $S \setminus \{r\}$ whose roots are the children of $r$ (the $v$ with $p(v) = r$); grouping the vertices by the root of their
tree, this is the same as a set partition of $S \setminus \{r\}$ with a rooted tree on each block. Conversely a root $r$ and such an assembly give a rooted tree, by declaring $r$ the parent of the
roots of the blocks. By Theorems~\ref{generating-functions:theorem-labelled-product} (with the class $\mathcal{X}$ for the root) and~\ref{generating-functions:theorem-exponential-formula}, the
exponential generating function $T = \sum_nt_nx^n/n!$ of rooted trees satisfies
\[
T = x\exp(T).
\]
By Theorem~\ref{generating-functions:theorem-lagrange} with $\phi = e^t$, $\frac{t_n}{n!} = [x^n]T = \frac{1}{n}[t^{n-1}]e^{nt} = \frac{1}{n}\cdot\frac{n^{n-1}}{(n - 1)!}$, so $t_n = n^{n-1}$. By
Lemma~\ref{generating-functions:lemma-trees-roots} there are $n^{n-1}/n = n^{n-2}$ trees.
\end{proof}

\begin{proof}[Second proof, by double counting (Pitman)]
Count the sequences of $n - 1$ steps that build a rooted tree on $[n]$ from the rooted forest with $n$ roots and no edges, where a step chooses a root $\rho$ of the current forest and a vertex $v$
not in the tree of $\rho$, and sets $p(\rho) = v$; the result is again a rooted forest, with one root fewer. If the current forest has $k$ roots, with trees of sizes $n_1, \ldots, n_k$, the number of
possible steps is $\sum_i(n - n_i) = kn - n = n(k - 1)$. So there are $\prod_{k=2}^{n}n(k - 1) = n^{n-1}(n - 1)!$ building sequences. On the other hand, each rooted tree is the result of exactly
$(n - 1)!$ of them, one for each order in which its $n - 1$ edges $\{v, p(v)\}$ are added: removing any set of parent assignments from a rooted tree leaves a rooted forest, so every order is a valid
building sequence, and the sequence determines the order. Hence there are $n^{n-1}$ rooted trees.
\end{proof}

\begin{remark}\label{generating-functions:remark-cayley}
The trees of Definition~\ref{generating-functions:definition-rooted-tree} are exactly the connected graphs without cycles on $S$, the trees of graph theory (Chapter~\ref{graph-theory}), where
Pr\"ufer's bijection between trees on $[n]$ and sequences in $[n]^{n-2}$ gives a third proof. For $n = 4$ there are $16$ trees: $12$ paths ($4!/2$ orderings of the vertices up to reversal) and $4$
stars (choose the centre). Lagrange inversion also counts rooted forests on $[n]$ with $k$ trees: $\binom{n - 1}{k - 1}n^{n-k}$ (Exercise~\ref{generating-functions:exercise-forests}).
\end{remark}

\section{Estimates}\label{generating-functions:section-estimates}

An exact formula is not always the most useful answer: one often wants to know how fast a sequence grows. This section proves elementary bounds, by comparing products term by term, for the central
binomial coefficients, the Catalan numbers and the Fibonacci numbers. They are the first instances of a general theory, sketched in the last remark, in which the growth of the coefficients is read
off the singularities of the generating function viewed as a complex function.

\begin{proposition}[Central binomial coefficients]\label{generating-functions:proposition-central-binomial-bounds}
For $n \geq 1$,
\[
\frac{4^n}{2n + 1} \leq \frac{4^n}{2\sqrt{n}} \leq \binom{2n}{n} \leq \frac{4^n}{\sqrt{2n + 1}} \leq 4^n.
\]
\end{proposition}

\begin{proof}
\emph{Step 1: a product formula.} Put $P_n = \binom{2n}{n}/4^n$. Since $\binom{2n}{n} = \frac{(2n)!}{n!^2}$ and $(2n)! = 2^nn! \cdot 1 \cdot 3 \cdots (2n - 1)$ (as in
Example~\ref{generating-functions:example-central-binomial}),
\[
P_n = \frac{1 \cdot 3 \cdots (2n - 1)}{2 \cdot 4 \cdots 2n} = \prod_{k=1}^{n}\frac{2k - 1}{2k}.
\]
\emph{Step 2: the upper bound.} For $k \geq 1$, $\frac{2k - 1}{2k} < \frac{2k}{2k + 1}$, since $(2k - 1)(2k + 1) = 4k^2 - 1 < 4k^2$. Hence
\[
P_n^2 = \prod_{k=1}^{n}\frac{2k - 1}{2k} \cdot \frac{2k - 1}{2k} < \prod_{k=1}^{n}\frac{2k - 1}{2k} \cdot \frac{2k}{2k + 1} = \prod_{k=1}^{n}\frac{2k - 1}{2k + 1} = \frac{1}{2n + 1},
\]
a telescoping product. So $P_n \leq 1/\sqrt{2n + 1}$.
\emph{Step 3: the lower bound.} For $k \geq 2$, $\frac{2k - 1}{2k} > \frac{2k - 2}{2k - 1}$, since $(2k - 1)^2 = 4k^2 - 4k + 1 > 4k^2 - 4k = 2k(2k - 2)$. Hence, keeping the factor $k = 1$ apart,
\[
P_n^2 \geq \frac14\prod_{k=2}^{n}\frac{2k - 1}{2k} \cdot \frac{2k - 2}{2k - 1} = \frac14\prod_{k=2}^{n}\frac{k - 1}{k} = \frac{1}{4n}.
\]
\emph{Step 4: the outer inequalities.} $2\sqrt{n} \leq 2n + 1$ because $4n \leq (2n + 1)^2$, and $\sqrt{2n + 1} \geq 1$.
\end{proof}

The outer, weaker bounds also have a one-line proof: $\binom{2n}{n}$ is the largest of the $2n + 1$ numbers $\binom{2n}{k}$ (Proposition~\ref{combinatorics:proposition-unimodal}), whose sum is
$4^n$. The inner bounds are much better: the ratio of the upper to the lower bound is $\frac{2\sqrt{n}}{\sqrt{2n + 1}} = \sqrt{\frac{4n}{2n + 1}} < \sqrt{2}$, so they determine $\binom{2n}{n}$ up to a
factor less than $\sqrt{2}$, for every $n$, while the outer bounds differ by the factor $2n + 1$.

\begin{corollary}\label{generating-functions:corollary-catalan-bounds}
For $n \geq 1$, $\displaystyle\frac{4^n}{2(n + 1)\sqrt{n}} \leq C_n \leq \frac{4^n}{(n + 1)\sqrt{2n + 1}}$. Moreover $\frac{C_{n+1}}{C_n} = \frac{2(2n + 1)}{n + 2}$, so the ratios increase strictly
from $1$ towards $4$, and $2^{n-1} \leq C_n \leq 4^n$ for all $n \geq 1$.
\end{corollary}

\begin{proof}
Divide the bounds of Proposition~\ref{generating-functions:proposition-central-binomial-bounds} by $n + 1$. From $C_n = \frac{(2n)!}{n!\,(n + 1)!}$,
\[
\frac{C_{n+1}}{C_n} = \frac{(2n + 2)(2n + 1)}{(n + 1)(n + 2)} = \frac{2(2n + 1)}{n + 2} = 4 - \frac{6}{n + 2},
\]
which increases strictly with $n$, equals $1$ for $n = 0$ and is at least $2$ for $n \geq 1$ (as $4n + 2 \geq 2n + 4$). So $C_1 = 1$ and $C_{n+1} \geq 2C_n$ give $C_n \geq 2^{n-1}$ by induction, and
$C_n \leq \binom{2n}{n} \leq 4^n$.
\end{proof}

\begin{counterexample}\label{generating-functions:counterexample-small-cases}
Bounds of this kind need their range checked separately. The tempting bound $C_n \geq 2^n$ is false for $n = 1, 2, 3, 4$ ($C_n = 1, 2, 5, 14$ against $2, 4, 8, 16$), and first holds at $n = 5$
($42 \geq 32$), after which it persists, since the ratio $C_{n+1}/C_n$ is then at least $2$. Likewise the lower bound $\frac{4^n}{2(n + 1)\sqrt{n}}$ is useless for small $n$ and good for large $n$,
while the bound $2^{n-1}$ is the opposite. An estimate is a statement for \emph{all} $n$ in a range, and the small cases are where it most often fails.
\end{counterexample}

\begin{proposition}[Fibonacci numbers]\label{generating-functions:proposition-fibonacci-bounds}
Let $\varphi = \frac{1 + \sqrt{5}}{2}$. For $n \geq 1$, $\varphi^{n-2} \leq F_n \leq \varphi^{n-1}$, and for $n \geq 0$, $F_n$ is the integer nearest to $\varphi^n/\sqrt{5}$.
\end{proposition}

\begin{proof}
\emph{The bounds}, by induction: they hold for $n = 1$ ($\varphi^{-1} \leq 1 \leq 1$) and $n = 2$ ($1 \leq 1 \leq \varphi$). If they hold for $n$ and $n + 1$, then, using $\varphi^2 = \varphi + 1$,
$F_{n+2} = F_{n+1} + F_n \leq \varphi^n + \varphi^{n-1} = \varphi^{n-1}(\varphi + 1) = \varphi^{n+1}$, and likewise $F_{n+2} \geq \varphi^{n-1} + \varphi^{n-2} = \varphi^n$.
\emph{Nearest integer.} By Binet's formula (Example~\ref{generating-functions:example-binet}), $F_n - \varphi^n/\sqrt{5} = -\psi^n/\sqrt{5}$ with $\psi = \frac{1 - \sqrt{5}}{2}$. Since $2 < \sqrt{5} < 3$,
$|\psi| = \frac{\sqrt{5} - 1}{2} < 1$, so $|\psi^n/\sqrt{5}| \leq \frac{1}{\sqrt{5}} < \frac{1}{2}$ for all $n \geq 0$.
\end{proof}

\begin{remark}[Analytic combinatorics]\label{generating-functions:remark-analytic}
The pattern behind these bounds is general. If the generating function $\sum_na_nx^n$ converges on a disc in $\CC$, the exponential growth rate $\limsup|a_n|^{1/n}$ is $1/\rho$, where $\rho$ is the radius
of convergence, and $\rho$ is the distance from $0$ to the nearest singularity of the function: $\frac{1}{4}$ for $\frac{1 - \sqrt{1 - 4x}}{2x}$, $\varphi^{-1}$ for $\frac{x}{1 - x - x^2}$. The
\emph{nature} of the singularity gives the subexponential factor: a square-root singularity $\sqrt{1 - x/\rho}$ produces $n^{-3/2}$, as for the Catalan numbers, and a simple pole produces a constant,
as for the Fibonacci numbers. The precise asymptotic formulas, such as
\[
\begin{aligned}
&\binom{2n}{n} \sim \frac{4^n}{\sqrt{\pi n}}, \qquad C_n \sim \frac{4^n}{\sqrt{\pi}\,n^{3/2}},\\
&\frac{D_n}{n!} \to \frac1e, \qquad p(n) \sim \frac{1}{4n\sqrt{3}}e^{\pi\sqrt{2n/3}},
\end{aligned}
\]
need real or complex analysis (Chapters~\ref{real-analysis} and~\ref{complex-analysis}); the last one is the Hardy--Ramanujan formula \cite{andrews-partitions}. The systematic theory is in
\cite{flajolet-sedgewick}. The first formula is consistent with Proposition~\ref{generating-functions:proposition-central-binomial-bounds}: $\frac{1}{2\sqrt{n}} < \frac{1}{\sqrt{\pi n}} < \frac{1}{\sqrt{2n + 1}}$
for $n \geq 1$, because $2n + 1 < \pi n < 4n$ when $n \geq 1$.
\end{remark}

\section{Exercises}\label{generating-functions:section-exercises}

\begin{exercise}\label{generating-functions:exercise-squares}
Show that $\sum_{n \geq 0}n^2x^n = \frac{x(1 + x)}{(1 - x)^3}$, and deduce $\sum_{k=0}^{n}k^2 = \frac{n(n + 1)(2n + 1)}{6}$ from Example~\ref{generating-functions:example-partial-sums}.
\end{exercise}

\begin{solution}
Let $E = xD$, which multiplies the coefficient of $x^n$ by $n$. From $\sum_nx^n = (1 - x)^{-1}$, $E(1 - x)^{-1} = x(1 - x)^{-2}$, and $E\bigl(x(1 - x)^{-2}\bigr) = x(1 - x)^{-2} + 2x^2(1 - x)^{-3} =
\frac{x(1 - x) + 2x^2}{(1 - x)^3} = \frac{x(1 + x)}{(1 - x)^3}$, by the Leibniz rule. Dividing by $1 - x$ gives the partial sums:
$\sum_n\bigl(\sum_{k \leq n}k^2\bigr)x^n = \frac{x + x^2}{(1 - x)^4}$, whose coefficient of $x^n$ is, by Example~\ref{generating-functions:example-geometric}(3),
$\binom{n + 2}{3} + \binom{n + 1}{3} = \frac{(n + 1)n}{6}\bigl((n + 2) + (n - 1)\bigr) = \frac{n(n + 1)(2n + 1)}{6}$.
\end{solution}

\begin{exercise}\label{generating-functions:exercise-compositions-one-two}
Show that the number of ways of writing $n$ as an ordered sum of parts equal to $1$ or $2$ is $F_{n+1}$, and that the number of ordered sums of parts $\geq 2$ is $F_{n-1}$ for $n \geq 1$.
\end{exercise}

\begin{solution}
A composition with $k$ parts in a set $A$ is a $k$-tuple, so by the rule of product the generating function is $\sum_k\bigl(\sum_{a \in A}x^a\bigr)^k = \frac{1}{1 - \sum_{a \in A}x^a}$. For $A = \{1, 2\}$
this is $\frac{1}{1 - x - x^2} = \frac{1}{x}\sum_nF_nx^n$, so the count is $F_{n+1}$. For $A = \{2, 3, \ldots\}$, $\sum_{a \geq 2}x^a = \frac{x^2}{1 - x}$ and
$\frac{1}{1 - x^2/(1 - x)} = \frac{1 - x}{1 - x - x^2} = 1 + \frac{x^2}{1 - x - x^2}$, whose coefficient of $x^n$ is $F_{n-1}$ for $n \geq 1$ (and $F_0 = 0$ for $n = 1$: there is no such composition of $1$).
\end{solution}

\begin{exercise}\label{generating-functions:exercise-central-convolution}
Prove by generating functions that
\[
\sum_{k=0}^{n}\binom{2k}{k}\binom{2n - 2k}{n - k} = 4^n,
\]
and show that $\sum_{k=0}^{n}C_kC_{n-k} = C_{n+1}$ is the same kind of identity.
\end{exercise}

\begin{solution}
The first identity is the square of $(1 - 4x)^{-1/2} = \sum_n\binom{2n}{n}x^n$, which is $(1 - 4x)^{-1}$ (Example~\ref{generating-functions:example-central-binomial}). For the second, $B = \sum_nC_nx^n$
satisfies $B = 1 + xB^2$ (Theorem~\ref{generating-functions:theorem-catalan}), so $B^2 = (B - 1)/x = \sum_nC_{n+1}x^n$: again the coefficients of a square of a series in closed form.
\end{solution}

\begin{exercise}\label{generating-functions:exercise-no-part-one}
Show that the number of partitions of $n \geq 1$ with no part equal to $1$ is $p(n) - p(n - 1)$, and that the number of partitions of $n$ in which no part occurs more than twice equals the number of
partitions of $n$ with no part divisible by $3$.
\end{exercise}

\begin{solution}
As in Theorem~\ref{generating-functions:theorem-partition-gf} and Example~\ref{generating-functions:example-partition-gf-variants}, the generating function of partitions with parts in a set
$A \subseteq \{1, 2, \ldots\}$ is $\prod_{i \in A}(1 - x^i)^{-1}$.
Removing the factor $i = 1$ multiplies $\sum_np(n)x^n$ by $1 - x$, which gives $p(n) - p(n - 1)$. For the second claim, a part $i$ occurring $0$, $1$ or $2$ times contributes the factor
$1 + x^i + x^{2i} = \frac{1 - x^{3i}}{1 - x^i}$; by Lemma~\ref{generating-functions:lemma-product-rules}(1) and (2),
\[
\prod_{i \geq 1}\frac{1 - x^{3i}}{1 - x^i} = \Bigl(\prod_{i \geq 1}\frac{1}{1 - x^i}\Bigr)\Bigl(\prod_{i \geq 1}\frac{1}{1 - x^{3i}}\Bigr)^{-1} = \prod_{3 \nmid i}\frac{1}{1 - x^i},
\]
the last step by splitting the index set (Lemma~\ref{generating-functions:lemma-product-rules}(3)). This generalises Glaisher's theorem (Corollary~\ref{generating-functions:corollary-glaisher}) with $3$ in place of $2$.
\end{solution}

\begin{exercise}\label{generating-functions:exercise-q-binomial-values}
Compute $\binom{4}{2}_q$, check that its value at $q = 1$ is $6$, and use Theorem~\ref{generating-functions:theorem-subspaces} to count the $2$-dimensional subspaces of $\mathbb{F}_2^4$.
\end{exercise}

\begin{solution}
$\binom{4}{2}_q = \frac{(1 - q^4)(1 - q^3)}{(1 - q)(1 - q^2)} = (1 + q^2)(1 + q + q^2) = 1 + q + 2q^2 + q^3 + q^4$; at $q = 1$ this is $6$, and the coefficients count the partitions in a $2 \times 2$
box by size: $\emptyset$; $(1)$; $(2), (1, 1)$; $(2, 1)$; $(2, 2)$. At $q = 2$ it is $1 + 2 + 8 + 8 + 16 = 35$ subspaces.
\end{solution}

\begin{exercise}\label{generating-functions:exercise-matchings}
A \emph{perfect matching} of a set is a partition into blocks of size $2$. Show with the exponential formula that the number of perfect matchings of $[2m]$ is $(2m - 1)(2m - 3) \cdots 3 \cdot 1$, and
that the number of involutions of $[n]$ is $\sum_k\frac{n!}{2^kk!\,(n - 2k)!}$.
\end{exercise}

\begin{solution}
Take $\mathcal{C}$ with one structure on each $2$-element set and none on other sets: $C = x^2/2$. By Theorem~\ref{generating-functions:theorem-exponential-formula} the exponential generating function is
$\exp(x^2/2) = \sum_m\frac{x^{2m}}{2^mm!}$, so the number is $\frac{(2m)!}{2^mm!} = 1 \cdot 3 \cdots (2m - 1)$ (Example~\ref{generating-functions:example-central-binomial}). For involutions,
$\exp(x + x^2/2) = e^x\exp(x^2/2)$, and by Theorem~\ref{generating-functions:theorem-labelled-product} the number is $\sum_k\binom{n}{2k}\frac{(2k)!}{2^kk!} = \sum_k\frac{n!}{2^kk!\,(n - 2k)!}$:
choose the $2k$ points that are moved, and a perfect matching of them.
\end{solution}

\begin{exercise}\label{generating-functions:exercise-ordered-partitions}
An \emph{ordered set partition} of $S$ is a set partition together with a linear order of its blocks. Show that their exponential generating function is $\frac{1}{2 - e^x}$, and compute the numbers
for $n \leq 4$. Show that the number of surjections $[n] \to [k]$ is $k!\,S(n, k)$.
\end{exercise}

\begin{solution}
An ordered set partition with $k$ blocks is a $k$-tuple of nonempty sets forming an ordered decomposition, so by Theorem~\ref{generating-functions:theorem-labelled-product} its generating function is
$(e^x - 1)^k$, and summing over $k$ (summable, as $\operatorname{ord}(e^x - 1)^k = k$) gives $\frac{1}{1 - (e^x - 1)} = \frac{1}{2 - e^x}$. Equivalently the number with $k$ blocks is $k!\,S(n, k)$
(Example~\ref{generating-functions:example-egf-partitions}), the $k!$ orders of the blocks. The values $\sum_kk!\,S(n, k)$ for $n = 0, \ldots, 4$ are $1, 1, 3, 13, 75$. A surjection $f \colon [n] \to [k]$ is
the same as the ordered set partition $(f^{-1}(1), \ldots, f^{-1}(k))$.
\end{solution}

\begin{exercise}\label{generating-functions:exercise-odd-cycles}
Show that the exponential generating function of the permutations all of whose cycles have odd length is $\sqrt{\frac{1 + x}{1 - x}}$, and that their number for $n = 2m$ is
$\bigl(1 \cdot 3 \cdots (2m - 1)\bigr)^2$.
\end{exercise}

\begin{solution}
By Example~\ref{generating-functions:example-egf-cycles}, the generating function is $\exp\bigl(\sum_{k \text{ odd}}x^k/k\bigr)$, and $\sum_{k \text{ odd}}\frac{x^k}{k} = \frac12\bigl(\log(1 + x) - \log(1 - x)\bigr)$
(Definition~\ref{generating-functions:definition-exp-log}). So the series is $(1 + x)^{1/2}(1 - x)^{-1/2}$ by Definition~\ref{generating-functions:definition-powers} and
Proposition~\ref{generating-functions:proposition-exp-log}(2). By Theorem~\ref{generating-functions:theorem-binomial-series}(2) this is $(1 + x)\bigl((1 + x)(1 - x)\bigr)^{-1/2} = (1 + x)(1 - x^2)^{-1/2}$,
and $(1 - x^2)^{-1/2} = \sum_m\binom{2m}{m}\frac{x^{2m}}{4^m}$ (Example~\ref{generating-functions:example-central-binomial} with $x^2/4$ in place of $x$). The coefficient of $x^{2m}$ is $\binom{2m}{m}/4^m$,
and $(2m)!\binom{2m}{m}/4^m = \Bigl(\frac{(2m)!}{2^mm!}\Bigr)^2 = \bigl(1 \cdot 3 \cdots (2m - 1)\bigr)^2$.
\end{solution}

\begin{exercise}\label{generating-functions:exercise-derangement-recurrence}
From $(1 - x)D(x) = e^{-x}$ (Example~\ref{generating-functions:example-labelled-product}), deduce $D_n = nD_{n-1} + (-1)^n$ for $n \geq 1$.
\end{exercise}

\begin{solution}
The coefficient of $x^n/n!$ in $xD(x)$ is $n!\frac{D_{n-1}}{(n - 1)!} = nD_{n-1}$. So comparing coefficients of $x^n/n!$ gives $D_n - nD_{n-1} = (-1)^n$.
\end{solution}

\begin{exercise}\label{generating-functions:exercise-sum-cubes-bernoulli}
Use Proposition~\ref{generating-functions:proposition-bernoulli}(3) to show $\sum_{k=0}^{m-1}k^3 = \bigl(\frac{m(m - 1)}{2}\bigr)^2$, and compute $B_8$.
\end{exercise}

\begin{solution}
With $n = 3$: $\frac14\bigl(B_0m^4 + 4B_1m^3 + 6B_2m^2 + 4B_3m\bigr) = \frac14\bigl(m^4 - 2m^3 + m^2\bigr) = \frac{m^2(m - 1)^2}{4}$. For $B_8$, the recurrence of
Proposition~\ref{generating-functions:proposition-bernoulli}(1) with $n = 8$, and $B_3 = B_5 = B_7 = 0$, reads
$B_0 + 9B_1 + 36B_2 + 126B_4 + 84B_6 + 9B_8 = 0$, that is $1 - \frac92 + 6 - \frac{21}{5} + 2 + 9B_8 = 0$, so $9B_8 = \frac{3}{10}$ and $B_8 = -\frac{1}{30}$.
\end{solution}

\begin{exercise}\label{generating-functions:exercise-ballot-numbers}
Show that the number of $k$-tuples of binary trees with $n$ nodes in total is $\frac{k}{2n + k}\binom{2n + k}{n}$, and check the cases $k = 1$ and $k = 2$.
\end{exercise}

\begin{solution}
The generating function is $B^k$ (rule of product), and $B^k = x^{-k}T^k$ with $T = xB$, $T = x(1 - T)^{-1}$ (Example~\ref{generating-functions:example-lagrange-catalan}). By
Theorem~\ref{generating-functions:theorem-lagrange} with $N = n + k$,
$[x^n]B^k = [x^N]T^k = \frac{k}{N}[t^{n}](1 - t)^{-N} = \frac{k}{n + k}\binom{2n + k - 1}{n}$ (Example~\ref{generating-functions:example-geometric}(3)), and
$\binom{2n + k}{n} = \frac{2n + k}{n + k}\binom{2n + k - 1}{n}$. For $k = 1$ this is $\frac{1}{2n + 1}\binom{2n + 1}{n} = \frac{1}{n + 1}\binom{2n}{n} = C_n$; for $k = 2$ it is
$\frac{1}{n + 1}\binom{2n + 2}{n} = C_{n+1}$, as in Exercise~\ref{generating-functions:exercise-central-convolution}.
\end{solution}

\begin{exercise}\label{generating-functions:exercise-forests}
Show that the number of rooted forests on $[n]$ with exactly $k$ trees is $\binom{n - 1}{k - 1}n^{n-k}$ for $1 \leq k \leq n$.
\end{exercise}

\begin{solution}
Grouping the vertices by the root of their tree, a rooted forest with $k$ trees is an assembly of $k$ rooted trees, so by Theorem~\ref{generating-functions:theorem-exponential-formula} the exponential
generating function is $T^k/k!$, with $T = x\exp(T)$ as in the proof of Theorem~\ref{generating-functions:theorem-cayley}. By Theorem~\ref{generating-functions:theorem-lagrange},
$[x^n]T^k = \frac{k}{n}[t^{n-k}]e^{nt} = \frac{k}{n}\cdot\frac{n^{n-k}}{(n - k)!}$, so the number is $\frac{n!}{k!}\cdot\frac{k}{n}\cdot\frac{n^{n-k}}{(n - k)!} = \frac{(n - 1)!}{(k - 1)!\,(n - k)!}n^{n-k}$.
For $k = 1$ this is Cayley's $n^{n-1}$; for $k = n$ it is $1$, the forest without edges.
\end{solution}

\begin{exercise}\label{generating-functions:exercise-motzkin}
A \emph{unary-binary tree} is a root together with an ordered list of $0$, $1$ or $2$ unary-binary trees. Show that the number with $n$ nodes is $\frac{1}{n}[t^{n-1}](1 + t + t^2)^n$, and compute it for
$n \leq 6$.
\end{exercise}

\begin{solution}
The generating function $T$ (counting nodes) satisfies $T = x(1 + T + T^2)$ by the rules of sum and product, and $T(0) = 0$, so Theorem~\ref{generating-functions:theorem-lagrange} applies with
$\phi = 1 + t + t^2$. The coefficient of $t^{n-1}$ in $(1 + t + t^2)^n$ is the sum of the multinomial coefficients $\binom{n}{a, b, c}$ with $b + 2c = n - 1$; for $n = 1, \ldots, 6$ the counts are
$1, 1, 2, 4, 9, 21$, the \emph{Motzkin numbers}. For example, for $n = 5$: $\frac15\bigl(\binom{5}{1, 4, 0} + \binom{5}{2, 2, 1} + \binom{5}{3, 0, 2}\bigr) = \frac{5 + 30 + 10}{5} = 9$.
\end{solution}

\begin{exercise}\label{generating-functions:exercise-ternary}
Compute the numbers of ternary trees with $n \leq 4$ nodes from the recurrence of Definition~\ref{generating-functions:definition-mary-tree}, and compare with
Corollary~\ref{generating-functions:corollary-fuss-catalan}.
\end{exercise}

\begin{solution}
With $t_n$ the number of ternary trees with $n$ nodes, $t_0 = 1$ and $t_n = \sum_{a+b+c=n-1}t_at_bt_c$. So $t_1 = t_0^3 = 1$; $t_2 = 3t_1t_0^2 = 3$ (the triples $(1, 0, 0)$ and its permutations);
$t_3 = 3t_2 + 3t_1^2 = 12$ (the $3$ permutations of $(2, 0, 0)$ and the $3$ of $(1, 1, 0)$); and $t_4 = 3t_3 + 6t_2t_1 + t_1^3 = 36 + 18 + 1 = 55$ (the $3$ permutations of $(3, 0, 0)$, the $6$ of
$(2, 1, 0)$, and $(1, 1, 1)$).
The formula $\frac{1}{2n + 1}\binom{3n}{n}$ gives $1, 1, 3, 12, 55$ for $n = 0, \ldots, 4$.
\end{solution}

\begin{exercise}\label{generating-functions:exercise-fibonacci-growth}
Show that $F_n \leq 2^{n-1}$ for $n \geq 1$ and that $F_{2n} \geq 2^{n-1}$ for $n \geq 1$, without using $\varphi$.
\end{exercise}

\begin{solution}
By induction: $F_1 = F_2 = 1$, and $F_{n+2} = F_{n+1} + F_n \leq 2^n + 2^{n-1} \leq 2^{n+1}$. For the second, $F_2 = 1$, $F_4 = 3 \geq 2$, and $F_{2n+2} = F_{2n+1} + F_{2n} \geq 2F_{2n}$, since
$F_{2n+1} \geq F_{2n}$.
\end{solution}

\begin{exercise}\label{generating-functions:exercise-not-rational}
Show that the generating function $\sum_n\binom{2n}{n}x^n$ is not a rational function: there are no polynomials $P, Q \in \QQ[x]$ with $Q(0) = 1$ and $\sum_n\binom{2n}{n}x^n = P/Q$.
\end{exercise}

\begin{solution}
If it were $P/Q$, its square $(1 - 4x)^{-1}$ would be $P^2/Q^2$, so $Q^2 = (1 - 4x)P^2$ in $\QQ[x]$. Compare the multiplicity of the irreducible factor $1 - 4x$ on both sides (unique factorisation
in $\QQ[x]$, which is Euclidean and hence a UFD by Theorem~\ref{rings:theorem-euclidean-pid-ufd}): it is even on the left and odd on the right. This contradiction shows that the sequence $\binom{2n}{n}$ satisfies no linear recurrence with constant
coefficients (Theorem~\ref{generating-functions:theorem-rational}).
\end{solution}
